\documentclass[12pt,a4paper]{report}

% -----------------------------
% Packages
% -----------------------------
\usepackage[utf8]{inputenc}
\usepackage[T1]{fontenc}
\usepackage{lmodern}
\usepackage{geometry}
\geometry{margin=1in}

\usepackage{amsmath,amssymb,amsfonts}
\usepackage{bm}
\usepackage{graphicx}
\usepackage{booktabs}
\usepackage{hyperref}
\usepackage{microtype}
\usepackage{enumitem}
\usepackage{listings}
\usepackage[T1]{fontenc}

\lstset{
  language=Python,
  basicstyle=\ttfamily\footnotesize,  % smaller mono font
  breaklines=true,                    % wrap long lines
  breakatwhitespace=true,             % wrap at spaces if possible
  columns=fullflexible,               % better use of horizontal space
  keepspaces=true,                    % respect spaces (indentation)
  showstringspaces=false,             % don't mark spaces in strings
  numberstyle=\tiny,
  stepnumber=1,
  frame=single,
}
\usepackage[utf8]{inputenc}  % if you're using pdfLaTeX
\usepackage[T1]{fontenc}
\usepackage{lmodern}

% Better looking hyperlinks
\hypersetup{
    colorlinks=true,
    linkcolor=blue,
    citecolor=blue,
    urlcolor=blue,
    pdftitle={ONTOΣ: A Volitional Ontology of Directedness, Coherence, and Drift},
    pdfauthor={Ryan O. van Gelder}
}

% For nicer math operators / commands
\newcommand{\Will}{W}
\newcommand{\Struct}{E}
\newcommand{\Act}{A}
\newcommand{\Espace}{\mathcal{E}(\Will)}
\newcommand{\Aspace}{\mathcal{A}(\Struct)}
\newcommand{\Ct}{C_t}
\newcommand{\Dt}{D_t}
\usepackage{newunicodechar}
\usepackage{amssymb} % for \blacktriangle etc.

% Map problematic Unicode chars to TeX equivalents
\newunicodechar{Δ}{\Delta}
\newunicodechar{▲}{\blacktriangle}

% For safety: if any of these still remain in the text,
% map them to simple ASCII approximations:
\newunicodechar{─}{-}
\newunicodechar{│}{|}
\newunicodechar{└}{+}
\newunicodechar{┘}{+}

% If you used Unicode arrows or ellipsis, you can add:
\newunicodechar{→}{\textrightarrow}
\newunicodechar{↓}{\ensuremath{\downarrow}}
\newunicodechar{…}{\dots}

% -----------------------------
% 0.1 Title Page
% -----------------------------
\title{\texorpdfstring{ONTO$\Sigma$: A Volitional Ontology of Directedness, Coherence, and Drift}
                     {ONTOS: A Volitional Ontology of Directedness, Coherence, and Drift}}


\author{\LARGE{In Legacy of Max Barzenkov}\\ \\
Formalized by Ryan O. van Gelder \\ \\
\small Institute of Mathematical Discovery}
\date{December 2025}

\begin{document}

\maketitle
\thispagestyle{empty}
\clearpage

% -----------------------------
% 0.2 Abstract / Executive Summary
% -----------------------------
\pagenumbering{roman}

\begin{abstract}
This report develops \emph{ONTO$\Sigma$}, a layered framework for understanding directedness, coherence, and drift across physical, cognitive, and social systems. At its core is the UTAM axiom,
\[
    \Will \to \Struct \to \Act,
\]
which asserts an ontological priority relation: \emph{will} (\Will) as a primitive operator of directedness, \emph{structure} (\Struct) as the manifestation or organization of will, and \emph{action} (\Act) as the realization of will through structure. In this reading, no structure exists without will, and no action is possible without a prior structure.

ONTO$\Sigma$ interprets will not as a psychological intention or divine decree, but as an impersonal, ontological vector of directionality that underlies all ordered becoming. Structures---from physical laws and organismic forms to neural networks and institutions---are then understood as locally stable organizations of this primordial directedness. Actions are concrete unfoldings of these structures in time. This vertical schema provides a metaphysical grounding for \emph{directedness} as such.

Building on this, the ONTO$\Sigma$ framework introduces two central dynamical notions: \emph{coherence} and \emph{drift}. Coherence, denoted \(\Ct \in [0,1]\), measures how well a system's current structure \(\Struct_t\) fits its incoming impulses \(I_t\) and its own actions \(\Act_t\). Drift, defined as \( \Dt = 1 - \Ct \), captures the loss of such fit as the environment changes or as the system fails to adapt. The \emph{Drift Law} states, qualitatively, that in a non-stationary environment any system whose internal structure does not adequately update in response to new impulses will inevitably lose coherence over time.

To move from metaphysics to dynamics, ONTO$\Sigma$ extends UTAM into a \emph{horizontal} layer. Time-indexed structures and actions obey:
\[
    \Act_t = F(\Struct_t, I_t), \qquad
    \Struct_{t+1} = G(\Struct_t, \Act_t, I_t, \Theta),
\]
where \(F\) is an action-generator and \(G\) a structure-updater governed by adaptation parameters \(\Theta\). These maps remain constrained by the vertical axiom: at every time \(t\), \(\Struct_t\) is still a manifestation of \Will, and \(\Act_t\) is still a realization of \Will\ through \(\Struct_t\).

In this report we propose two minimal algorithmic schemes that operationalize \(F\) and \(G\) in adaptive systems:

\begin{itemize}[leftmargin=2em]
    \item The \textbf{I$^2$C cycle} (Impulse $\to$ Interpretation $\to$ Coherence) models how a system maps an impulse \(I_t\) into an internal representation \(Z_t\) and selects an action \(\Act_t\) that maximizes a coherence functional \(C(Z_t, \Act, \Struct_t)\).
    \item The \textbf{$\Delta \Struct$ update} (\(\Delta E\)) encodes how the system updates its structure \(\Struct_t\) to preserve or increase coherence under drift, for example via a gradient-based rule
    \[
        \Struct_{t+1} = \Struct_t + \eta \, \nabla_{\Struct} C(\Struct_t, I_t, \Act_t),
    \]
    where \(\eta\) is a learning rate. More generally, any update rule that tends to increase coherence \emph{in expectation} counts as a $\Delta \Struct$ process.
\end{itemize}

We then introduce the notion of a \emph{volitional phase space}, in which the space of possible structures \(\mathcal{E}(\Will)\) is treated as a manifold equipped with a coherence landscape. Trajectories of \(\Struct_t\) under I$^2$C + $\Delta \Struct$ are interpreted as gradient flows on this landscape, while environmental drift deforms the landscape itself. Psychological states such as flow, burnout, anxiety, and aging are interpreted as dynamical regimes or attractor structures within this volitional phase space. Meta-structural dynamics---arising through consciousness, culture, and technology---correspond to reshaping and expanding the accessible region of \(\mathcal{E}(\Will)\) without altering \Will\ itself.

The second half of the report outlines a research program that connects ONTO$\Sigma$ to empirical and technical domains. In neuroscience, we propose relating coherence \(C_t\) to patterns of neural synchronization and functional connectivity, and interpreting drift in terms of maladaptive plasticity. In AI and robotics, we discuss how $\Delta \Struct$ controllers can yield systems more robust to distribution shift than classical setpoint- or reward-based approaches. In psychology, we interpret meaning, selfhood, and burnout as regimes of coherence and drift, suggesting new ways to structure therapeutic interventions as $\Delta \Struct$ training. At the level of social and institutional systems, we sketch how civilizational coherence and collapse may be understood as large-scale UTAM dynamics.

Throughout, ONTO$\Sigma$ remains an open, layered framework. It does not seek to replace domain-specific theories, but to provide a unified ontological, dynamical, and geometric vocabulary for directedness, coherence, and drift. The central empirical question it poses is whether will (\Will), understood as an ontological operator of directedness, yields deeper insight into adaptive phenomena than can be obtained through purely mechanistic or purely goal-centric models. The report concludes with executable toy models and proposals for cross-disciplinary experiments designed to test and refine this claim.
\end{abstract}

\clearpage

% -----------------------------
% 0.3 Table of Contents
% -----------------------------
\tableofcontents
\clearpage

% -----------------------------
% 0.4 Notation and Conventions
% -----------------------------
\section*{0.4 Notation and Conventions}
\addcontentsline{toc}{section}{0.4 Notation and Conventions}

\subsection*{0.4.1 Symbols}

Table \ref{tab:notation} summarizes the main symbols used throughout this report.

\begin{table}[h!]
    \centering
    \begin{tabular}{ll}
        \toprule
        \textbf{Symbol} & \textbf{Meaning} \\
        \midrule
        $\Will$ & Primitive ontological operator of directedness (will) \\
        $\mathcal{E}(\Will)$ & Space of possible structures that can manifest $\Will$ \\
        $E_t \in \mathcal{E}(\Will)$ & Concrete structure at time $t$ \\
        $\mathcal{A}(E)$ & Space of possible actions realizable through structure $E$ \\
        $A_t \in \mathcal{A}(E_t)$ & Action at time $t$ \\
        $I_t$ & Environmental impulse at time $t$ \\
        $C_t \in [0,1]$ & Coherence measure at time $t$ \\
        $D_t = 1 - C_t$ & Drift measure at time $t$ \\
        $F$ & Action-generator implementing the I$^2$C cycle \\
        $G$ & Structure-updater implementing the $\Delta E$ (Delta-$E$) principle \\
        $\eta$ & Learning rate / adaptation speed \\
        $\Theta$ & Adaptation parameters (e.g.\ plasticity, meta-learning rates) \\
        \bottomrule
    \end{tabular}
    \caption{Main symbols and their meanings.}
    \label{tab:notation}
\end{table}

\subsection*{0.4.2 Typographical Conventions}

We adopt the following typographical conventions throughout the document:

\begin{itemize}[leftmargin=2em]
    \item \emph{Concepts} (e.g.\ \emph{coherence}, \emph{drift}, \emph{directedness}) are set in italics on first introduction.
    \item \textbf{Key definitions} and principles (e.g.\ the \textbf{UTAM axiom}, the \textbf{Drift Law}) are typeset in bold when formally stated.
    \item Mathematical objects such as functions, sets, and spaces are written in standard math fonts (e.g.\ $F$, $G$, $\mathcal{E}(\Will)$).
    \item The notation I$^2$C is used for the \emph{Impulse $\to$ Interpretation $\to$ Coherence} cycle, and $\Delta E$ (or $\Delta \Struct$) for the structure-update principle.
    \item When referring to the ONTO$\Sigma$ series of texts (ONTO$\Sigma$ I, II, IIb, III), we use roman numerals in small caps where appropriate.
\end{itemize}

\clearpage

% Switch to arabic numbering for main body
\pagenumbering{arabic}

% ================================================
% Part I – Foundations of the ONTOΣ Framework
% ================================================

\part{Foundations of the ONTO$\Sigma$ Framework}

\chapter{Historical and Conceptual Background}

\section{Motivations from \emph{Synthetic Conscience}}

ONTO$\Sigma$ did not arise in a vacuum. It emerged as the closing movement of the broader \emph{Synthetic Conscience} sequence, which began as an attempt to understand and engineer adaptive, self-regulating systems---artificial agents, cognitive architectures, and hybrid human--machine setups---that could maintain internal stability in changing environments.

In early work, the focus was almost exclusively on \emph{mechanisms}: feedback loops, controllers, and learning algorithms that seemed to produce robust behaviour. Certain patterns kept reappearing across very different domains:

\begin{itemize}[leftmargin=2em]
    \item \textbf{$\Delta E$-like updates}, where some internal structure was continually adjusted in response to discrepancies between expected and observed input.
    \item \textbf{Thermostat dynamics}, where systems maintained key variables within a viable range by counteracting deviations---a minimal model of self-regulation.
    \item \textbf{Coherence functions}, which served as informal metrics for how well an internal model, policy, or narrative fit the actual demands of its environment.
\end{itemize}

These motifs appeared in neural-network training, adaptive control schemes, and even in the phenomenology of personal change. In each case there was some ``structure''---weights, parameters, beliefs, habits---that had to be preserved or updated in order for the system to remain viable. The language of \emph{coherence} and \emph{drift} suggested itself naturally: certain configurations of the system were well-matched to the world; others gradually lost their fit as conditions changed.

At first, this seemed like an engineering story: design better controllers, craft more flexible learning rules, tune hyperparameters. But the recurrence of the same abstract pattern across such diverse levels (from low-level thermostats to high-level narrative self-understanding) pointed to something deeper. It was not just that systems \emph{happened} to maintain coherence; they behaved as if there were a fundamental \emph{tendency} toward structured, directed adaptation.

This is where ``mere engineering'' began to show its limits. Standard tools could describe \emph{how} to maintain a setpoint or update a policy, but they offered no satisfying account of \emph{why} systems exhibit this stubborn drive to persist, adapt, and reorganize themselves---why there is directedness at all, rather than indifferent chaos. ONTO$\Sigma$ is the attempt to name and formalize that underlying layer: to treat directedness itself not as an accident of particular mechanisms, but as an ontological feature of reality. The UTAM axiom,
\[
    W \to E \to A,
\]
is the concise expression of this move: behind every adaptive mechanism there is will (\Will) as an operator of directedness, structure (\Struct) as its organization, and action (\Act) as its unfolding.

\section{Classical Philosophical Influences}

Although ONTO$\Sigma$ grew out of concrete work in adaptive systems and Synthetic Conscience, it consciously situates itself within a longer philosophical tradition that has wrestled with the inner character of reality.

A first obvious reference is Arthur Schopenhauer. For Schopenhauer, the \emph{world as representation} (the ordered manifold of objects in space and time) is grounded in an underlying \emph{world-will}, an insatiable, blind striving that constitutes the thing-in-itself. All phenomena are, on his view, objectifications of this will. ONTO$\Sigma$ shares with Schopenhauer the intuition that behind appearances there is a unitary, non-representational core---something more fundamental than particular laws or configurations. However, it diverges in a crucial respect: where Schopenhauer's will is characteristically blind and goalless, ONTO$\Sigma$ treats will (\Will) as a primitive \emph{operator of directedness} that is inherently structured and structurable. Will is not a chaotic surge that happens to produce order; it is that in virtue of which order and directedness are possible at all.

Henri Bergson's notion of \emph{élan vital}---a vital impetus that drives creative evolution---also resonates with ONTO$\Sigma$. Bergson emphasizes duration, novelty, and the irreducible creativity of life. In ONTO$\Sigma$, will is likewise not a static principle but something that unfolds through increasingly complex structures, giving rise to novelty and expanded capacities. Yet, whereas Bergson's account remains largely at the level of life and consciousness,


\chapter{The UTAM Axiom: Will, Structure, Action}

The central proposal of ONTO$\Sigma$ is that all directed phenomena---from physical processes and biological adaptation to cognition, social organization, and meaning---can be understood as manifestations of a single underlying pattern. This pattern is captured by the \emph{UTAM axiom}, which posits a strict ontological ordering:
\[
    \Will \to \Struct \to \Act.
\]
In this chapter we unpack this schema, clarifying what is meant by will, structure, and action, and in what precise sense the arrows express \emph{priority} rather than temporal sequence.

\section{The Core Schema \(W \to E \to A\)}

We begin by stating UTAM in the form of three interconnected postulates.

\begin{description}[leftmargin=2em]
    \item[Primacy of Will (\(\Will\)).] There exists a fundamental ontological operator of directedness, called \emph{will}, which is prior to any particular structure or event. Will is that in virtue of which reality can exhibit ordered tendencies, trajectories, or orientations at all.
    \item[Mediation by Structure (\(\Struct\)).] Will does not appear ``naked'' in the world. It is always realized through \emph{structures}: organized configurations of matter, energy, information, or social order that channel and constrain its directedness.
    \item[Actualization as Action (\(\Act\)).] Every concrete action or event is the unfolding of some structure over time, and thus an indirect manifestation of will through that structure.
\end{description}

The schema
\[
    \Will \to \Struct \to \Act
\]
should be read informally as: \emph{\Will{} provides directedness, \Struct{} embodies it, \Act{} enacts it}. This is the core of UTAM: all directed behaviour, from the simplest physical relaxation process to the most complex human project, is understood as will taking form as a structure and then unfolding as an action.

\section{Will (\(W\)) as Ontological Operator of Directedness}

The notion of \emph{will} used in ONTO$\Sigma$ is not the everyday notion of wanting, intending, or choosing. It is not a psychological state belonging to an individual subject, nor a divine fiat imposed from outside the world. Rather, \Will{} is posited as a primitive \emph{ontological operator of directedness}: a ``vector of being'' that makes it possible for reality to exhibit ordered trends, biases, and tendencies instead of mere undifferentiated chaos.

It is helpful to distinguish \Will{} from several nearby notions:

\begin{itemize}[leftmargin=2em]
    \item \textbf{Subjective will.} When we speak of a person's will---their desires, decisions, or intentions---we are already presupposing a highly structured organism, a nervous system, a language, a history. These are \emph{manifestations} of \Will{}, not its ground. In UTAM, individual willing is a particular, derivative form of the more basic ontological directedness.
    \item \textbf{Divine will.} Theological traditions often speak of a divine will that orders the universe according to a plan. ONTO$\Sigma$ does not presuppose or deny such a notion; it simply works at a different level. \Will{} here is not a personal agent with preferences, but a structural condition of possibility for any order or purposiveness whatsoever.
    \item \textbf{Physical forces and fields.} Modern physics describes interactions in terms of forces, fields, and symmetries. These are highly specific structures, encoded in mathematical laws. UTAM treats such laws as elements of \(\mathcal{E}(\Will)\), the space of possible structures, rather than as ultimate primitives. Forces are \emph{shaped will}, not substitutes for it.
\end{itemize}

There is an analogy---and a contrast---to the Kantian and Schopenhauerian idea of the \emph{thing-in-itself}. Just as the thing-in-itself is posited as the hidden reality behind appearances, \Will{} is posited as the hidden operator behind all structured directedness. But unlike an entirely opaque noumenon, \Will{} is characterized positively as \emph{directional}: it is that which is expressed wherever there is a stable tendency, a bias in state space, a gradient that systems ``follow'' in their evolution.

In the UTAM framework, then, \Will{} is not a mysterious extra substance added to the world, but a name for the fact that the world has an intrinsic propensity to organize, persist, and develop along certain lines rather than others. It is the ontological root of directedness.

\section{Structure (\(E\)) as Structured Will}

If \Will{} is the primitive operator of directedness, \emph{structure} \(\Struct\) is the way this directedness becomes concrete. Structures are the organized patterns, laws, and configurations through which will is channelled and expressed.

Examples of structures include:

\begin{itemize}[leftmargin=2em]
    \item physical laws and boundary conditions in physics;
    \item morphologies and regulatory networks in biology;
    \item neural architectures and synaptic weight patterns in brains and artificial networks;
    \item grammars, concepts, and schemas in cognition;
    \item norms, institutions, and infrastructures in social systems.
\end{itemize}

All of these can be seen as elements of \(\mathcal{E}(\Will)\): the space of possible structures that can manifest will. In this sense, \Struct{} is not an independent substance that exists alongside \Will{}, but a \emph{local form of will}. To say that a particular structure exists is to say that will has taken on a specific, organized configuration that constrains and shapes how it can be enacted.

Two points are crucial here:

\begin{enumerate}[leftmargin=2em]
    \item Structures are \emph{embodiments of directedness}. They encode which trajectories are possible, likely, or stable, and which are not. A valley structure in a landscape channels water; a regulatory network channels developmental processes; a policy network channels behaviour.
    \item Structures are themselves \emph{products of prior processes}. They arise through the co-evolution of systems and environments. In later parts of this report we describe how \(\Struct_t\) is continually updated by $\Delta \Struct$ dynamics in response to impulses \(I_t\), but even at the purely ontological level, it is important to see \Struct{} as the crystallized history of will’s prior unfoldings.
\end{enumerate}

In short, \Struct{} is \emph{structured will}: the relatively stable, organized aspect of reality through which directedness can be expressed in a consistent way.

\section{Action (\(A\)) as Will-Through-Structure}

\emph{Action} \(\Act\) is the final term in the UTAM schema: the concrete unfolding of will through structure in time. Given a structure \(\Struct\) that channels and constrains directedness, actions are the actual state transitions, behaviours, or events that occur when this structured will is brought into contact with impulses and conditions.

At different levels, this may look like:

\begin{itemize}[leftmargin=2em]
    \item a physical system relaxing along a gradient, following the constraints encoded in a potential function;
    \item an organism moving, feeding, reproducing in ways organized by its anatomy and neural circuitry;
    \item a cognitive agent making a decision, drawing an inference, or uttering a sentence in line with its internal model and goals;
    \item an institution enacting a policy, a market reallocating resources, a legal system applying a statute.
\end{itemize}

In each case, we can say: \Act{} is \emph{will-through-structure}. It is not a free-floating event, but the manifestation of \Will{} operating under the specific constraints and possibilities encoded in \Struct. Every action presupposes a structure, because without some organization, there is no sense in which one trajectory rather than another counts as a realization of directedness.

This perspective also inverts a common intuition: instead of thinking that structures are built out of repeated actions, ONTO$\Sigma$ treats repeated actions as the ongoing enactment of underlying structures, which themselves are structured will. Of course, historically and dynamically structures are shaped by actions (as we will see in the horizontal UTAM), but ontologically the dependence runs from \Will{} to \Struct{} to \Act.

\section{The UTAM Priority Relation}

We can now state the \textbf{UTAM axiom} in its concise form:
\[
    \Will \to \Struct \to \Act.
\]

This arrow notation expresses a strict \emph{priority relation}:

\begin{itemize}[leftmargin=2em]
    \item There is \textbf{no} structure without will: any \(\Struct\) is a structuring of \Will{}, not an independent substrate.
    \item There is \textbf{no} action without structure: any \(\Act\) is the unfolding of some \(\Struct\), not an unconditioned flash out of nowhere.
\end{itemize}

It is essential to emphasize that, at this stage, the arrows \(\to\) are \emph{ontological}, not yet \emph{dynamical}. They do \emph{not} mean that \Will{} first occurs, then later produces \Struct{}, which then later produces \Act{} in linear time. Rather, they encode dependence:

\begin{itemize}[leftmargin=2em]
    \item Whenever you find a structure, you can say: there is will, organized in this way.
    \item Whenever you find an action, you can say: there is will, organized in this way, unfolding like this.
\end{itemize}

In later chapters we will introduce explicit dynamics and feedback loops---the horizontal UTAM---in which structures and actions evolve over time:
\[
    \Act_t = F(\Struct_t, I_t), \qquad
    \Struct_{t+1} = G(\Struct_t, \Act_t, I_t, \Theta).
\]
Those equations \emph{presuppose} the UTAM priority relation. They tell us how particular \(\Struct_t\) and \(\Act_t\) change, but not why there are structured, directed processes at all. The UTAM axiom is the answer to that deeper question: directedness originates in will, takes form as structure, and becomes visible as action.

\chapter{Directedness, Coherence, and Drift}

The UTAM axiom,
\[
    \Will \to \Struct \to \Act,
\]
tells us how will, structure, and action are related at the level of ontological priority. But to understand adaptive systems as they actually unfold in time, we need a more fine-grained vocabulary. In particular, we need to distinguish between mere \emph{causal happening} and genuinely \emph{directed} processes, and we need concepts that capture when a system is well-aligned with its world and when it is not. This chapter introduces three such concepts: \emph{directedness}, \emph{coherence}, and \emph{drift}.

\section{Directedness vs.\ Mere Causality}

Physical theories, especially in their classical formulations, are often read as describing a world in which ``things just happen'' according to timeless laws. Given an initial condition and a dynamical equation, the future state is determined (or stochastically constrained). In such a view, talk of things \emph{going somewhere}---of tendencies, aims, or directions---is often treated as a mere shorthand for underlying causal mechanisms.

ONTO$\Sigma$ insists on a distinction between \emph{causality} and \emph{directedness}:

\begin{itemize}[leftmargin=2em]
    \item \textbf{Mere causality} refers to ordered regularities of the form ``if a system is in state $x$ at time $t$, then with high probability it will be in state $y$ at time $t + \Delta t$.'' This is a description of how states succeed one another, but it is neutral as to whether the succession has any intrinsic orientation or \emph{preference}.
    \item \textbf{Directedness} refers to the presence of a \emph{structured bias} in the evolution of states: some trajectories are not just possible but \emph{favoured}, in the sense that the system tends to move toward them, maintain them, or return to them under perturbation.
\end{itemize}

In state-space terms, directedness shows up as gradients, attractors, and invariant sets that are not arbitrary but systematically related to the system's organisation. A stone sliding down a potential well, an organism seeking nutrients, a controller steering a process back to setpoint, a person pursuing a life project---all of these are instances where the dynamics are not indifferent: there is a sense in which the system ``goes somewhere'' rather than simply evolving.

In the UTAM framework, this ``going somewhere'' is precisely where \Will{} enters. Causal laws tell us \emph{how} states change; \Will{} is what makes it intelligible that there are stable directions, preferred regions, and robust attractors at all. Directedness is thus understood as a manifestation of will through structure: the world is not only causally ordered, it is \emph{oriented}.

\section{Coherence as Alignment of \(W\)–\(E\)–\(A\)}

If directedness is the general tendency of systems to ``go somewhere'' rather than nowhere in particular, \emph{coherence} is a measure of \emph{how well} a particular system is managing this directedness at a given moment. Intuitively, a system is coherent when:

\begin{itemize}[leftmargin=2em]
    \item its current structure \(\Struct_t\) is well-suited to the impulses \(I_t\) it is receiving,
    \item its actions \(\Act_t\) follow intelligibly from that structure,
    \item and the resulting unfolding does not undermine the very organisation that makes the system what it is.
\end{itemize}

In the language of UTAM, we can say that there is \emph{coherence between \(\Will, \Struct_t, \Act_t\) at time \(t\)} when will's directedness, the system's current form, and its actual behaviour line up in a mutually reinforcing way. We will later represent this by a coherence measure \(\Ct \in [0,1]\), but even before formalization the idea is clear:


\begin{itemize}[leftmargin=2em]
    \item At \textbf{high coherence}, the system feels and appears ``in tune'' with its world. Its responses are appropriate, its internal organisation is not being eroded by its own actions, and small deviations can be corrected.
    \item At \textbf{low coherence}, the system is out of step. Its structure no longer fits the demands placed upon it, its actions may be self-sabotaging, and its internal organisation is under strain.
\end{itemize}

We can distinguish between:

\begin{description}[leftmargin=2em]
    \item[Local coherence.] The alignment of \(\Will\)–\(\Struct_t\)–\(\Act_t\) at a particular time, in a particular context. A controller might be locally coherent with respect to one operating condition and incoherent with respect to another.
    \item[Long-term coherence regimes.] Stable patterns in which a system maintains high coherence over extended periods and across changing conditions. These regimes are central to ONTO$\Sigma$'s account of \emph{self}, \emph{meaning}, and \emph{integrity}: a person or institution ``has integrity'' insofar as it sustains a coherent way of being across the drifts of circumstance.
\end{description}

From the first-person perspective, high coherence often corresponds to a sense of meaningful engagement, flow, or authenticity, while low coherence corresponds to confusion, alienation, or fragmentation. ONTO$\Sigma$ will later use this connection to reinterpret these phenomenological states as regimes in the system's volitional phase space.

\section{Drift as Loss of Coherence}

\emph{Drift} is the process by which coherence is lost over time. Formally, we will define drift as
\[
    D_t = 1 - C_t,
\]
but conceptually it is the name for what happens when a structure \(\Struct_t\) no longer fits the impulses \(I_t\) it receives or the actions \(\Act_t\) it takes.

Drift can arise in several ways:

\begin{itemize}[leftmargin=2em]
    \item The environment changes (the distribution of \(I_t\) shifts), but the system's structure \(\Struct_t\) does not update accordingly.
    \item The system's own actions \(\Act_t\) gradually erode its internal organisation (self-undermining behaviour).
    \item Internal processes (e.g.\ noise, uncorrected errors) progressively distort \(\Struct_t\) even in a stable environment.
\end{itemize}

These possibilities are familiar across many domains:

\begin{itemize}[leftmargin=2em]
    \item In \textbf{machine learning}, \emph{concept drift} occurs when the data distribution changes relative to the training regime. A model that once performed well begins to fail, not because its internal structure has been corrupted, but because that structure no longer matches the actual $I_t$ it encounters.
    \item In \textbf{control systems}, controllers tuned for one operating regime may become unstable or ineffective when plant dynamics or loads change. The system drifts out of the region where its control law yields coherent behaviour.
    \item In \textbf{individual psychology}, burnout and depersonalization can be seen as cases where a person's internal structures (beliefs, habits, commitments) become increasingly misaligned with the demands placed upon them and with their own deeper needs. Actions that once made sense no longer feel connected to a coherent whole.
    \item In \textbf{institutions and civilizations}, laws, norms, and infrastructures that once maintained social coherence may lag behind technological, ecological, or cultural shifts. Over time, the gap between structural form and actual conditions widens, leading to dysfunction or collapse.
\end{itemize}

In each case, what is at stake is the same: a loss of alignment between will's directedness, the structures through which it is expressed, and the actions that actually occur. Drift is thus not merely error or failure in a specific task; it is a deeper misfit between a system's way of being and the world it inhabits.

Later, when we introduce the Drift Law, we will see that drift is not an accident but an \emph{inevitable consequence} of non-stationary environments combined with insufficient structural adaptation. ONTO$\Sigma$'s dynamic layer (I$^2$C and $\Delta \Struct$) is precisely about understanding how coherence can be maintained---or at least gracefully degraded---under such conditions.

\section{UTAM as a Metaphysics of Directedness}

We are now in a position to state more clearly what is distinctive about ONTO$\Sigma$ as a philosophical proposal. It is not merely a collection of useful metaphors or engineering heuristics. It is a \emph{metaphysics of directedness}.

The key claim is that \Will{} is not just a way of talking about certain regularities; it is an \emph{ontological primitive}. Alternative metaphysical options include:

\begin{itemize}[leftmargin=2em]
    \item \textbf{Law-based views}, where fundamental reality consists of inert entities governed by timeless laws. Directedness is then reduced to particular initial conditions and the form of those laws.
    \item \textbf{Powers or dispositions}, where entities are endowed with intrinsic powers that tend to produce certain effects. Directedness becomes a property of these powers, but the overall ontology may remain fragmented.
    \item \textbf{Information-centric views}, where the world is fundamentally informational and dynamics are transformations of information. Directedness is then an emergent feature of informational flows or optimization principles.
\end{itemize}

ONTO$\Sigma$ does not deny the usefulness of these perspectives. Laws, powers, and information are indispensable tools for modelling. But it argues that they leave something unexplained: the pervasive fact that systems not only change, but \emph{organize} themselves, maintain coherence, and respond to drift in ways that look like the expression of a deeper tendency.

By positing \Will{} as an ontological operator of directedness, UTAM provides a unifying backdrop:

\begin{itemize}[leftmargin=2em]
    \item Laws, powers, and informational structures are reinterpreted as elements of \(\mathcal{E}(\Will)\), specific ways in which will is organized.
    \item Directed processes in physics, biology, cognition, and society are all manifestations of \Will{} unfolding through different \(\Struct\) to produce different \(\Act\).
    \item Coherence and drift become cross-cutting concepts: they apply wherever \Will{} is structured and enacted, independent of the particular domain.
\end{itemize}

In this sense, UTAM is offered as a \emph{minimal enrichment} of our ontology: instead of multiplying substances or invoking ad hoc teleologies, it introduces a single, general principle---will as directedness---and shows how it can ground and unify the patterns already observed in diverse sciences and practices. The subsequent parts of this report are devoted to making this claim precise, dynamic, and empirically fruitful.


% ================================================
% Part II – Dynamics: From Ontology to Time
% ================================================

\part{Dynamics: From Ontology to Time}

\chapter{Vertical and Horizontal UTAM}

In Part I we introduced the UTAM axiom,
\[
    \Will \to \Struct \to \Act,
\]
as a statement about \emph{ontological priority}: will as primitive, structure as its manifestation, action as its realization. We now extend this static picture into the temporal domain, distinguishing a \emph{vertical} layer of dependence from a \emph{horizontal} layer of dynamics and feedback.

\section{Vertical UTAM: Ontological Priority}

We begin by restating the vertical layer in more formal terms.

\subsection*{Primitive will and spaces of manifestation}

\begin{itemize}[leftmargin=2em]
    \item \(\Will\) denotes the \textbf{primitive ontological operator of directedness}. It is not a state or an event, but the underlying source of directionality in reality.
    \item \(\mathcal{E}(\Will)\) denotes the \textbf{space of possible structures} that can manifest \(\Will\). Each element \(E \in \mathcal{E}(\Will)\) is a way in which will is organized: a law, a morphology, a model, a norm, and so on.
    \item For each structure \(E\), \(\mathcal{A}(E)\) denotes the \textbf{space of possible actions} realizable through that structure. Each \(A \in \mathcal{A}(E)\) is a possible unfolding of \(E\) in time.
\end{itemize}

Formally, we can introduce two abstract selection maps:
\[
    \Phi : \Will \leadsto E \in \mathcal{E}(\Will),
    \qquad
    \Psi : E \leadsto A \in \mathcal{A}(E).
\]

These maps should \emph{not} be read as temporal transitions. They encode \emph{ontological dependence}:

\begin{itemize}[leftmargin=2em]
    \item \(\Phi\) expresses that any concrete structure \(E\) is a structuring of \(\Will\): there is no structure that is not an organization of will.
    \item \(\Psi\) expresses that any concrete action \(A\) is a realization of \(\Will\) through some structure \(E\): there is no action that is not the unfolding of some organization.
\end{itemize}

Thus, at the vertical level, UTAM asserts:
\[
    \text{No } E \text{ without } \Will, \qquad
    \text{No } A \text{ without } E.
\]

This remains entirely atemporal. It says nothing yet about how particular \(E\) and \(A\) change, only about how they \emph{stand} in dependence on \(\Will\).

\section{Horizontal UTAM: Temporal Feedback and Adaptation}

To describe actual systems as they exist in time, we introduce an explicit temporal index and move to a \emph{horizontal} layer.

Let:

\begin{itemize}[leftmargin=2em]
    \item \(E_t \in \mathcal{E}(\Will)\) be the effective structure of the system at time \(t\),
    \item \(I_t\) be the \textbf{impulses} impinging on the system from its environment at time \(t\),
    \item \(A_t \in \mathcal{A}(E_t)\) be the action taken at time \(t\),
    \item \(\Theta\) be a collection of \textbf{adaptation parameters} (e.g.\ plasticity, learning rates, meta-parameters).
\end{itemize}

At the horizontal, dynamical level we describe the system by two core maps:
\begin{equation}
\begin{aligned}
    A_t &= F(E_t, I_t), \\
    E_{t+1} &= G(E_t, A_t, I_t, \Theta).
\end{aligned}
\label{eq:horizontal-utam}
\end{equation}

Here:

\begin{itemize}[leftmargin=2em]
    \item \(F\) is the \textbf{action-generator}: given the current structure \(E_t\) and impulse \(I_t\), it produces an action \(A_t\). In later chapters we will refine \(F\) as the I$^2$C (Impulse $\to$ Interpretation $\to$ Coherence) cycle.
    \item \(G\) is the \textbf{structure-updater}: given the current structure \(E_t\), action \(A_t\), impulse \(I_t\), and adaptation parameters \(\Theta\), it produces the next structure \(E_{t+1}\). This will later be instantiated as the $\Delta E$ (Delta-\(E\)) principle.
\end{itemize}

Equations \eqref{eq:horizontal-utam} describe how particular structures and actions evolve in time, under the influence of environmental impulses and internal adaptation. They sit ``horizontally'' on top of the vertical UTAM, which specifies what kinds of entities \(E_t\) and \(A_t\) are.

\section{Closed-Loop Structure: Vertical and Horizontal Together}

Putting the vertical and horizontal layers together, we obtain the following picture:

\begin{center}
\begin{minipage}{0.9\textwidth}
\begin{verbatim}
      (vertical, ontological)
          W
          |
         E_t
          |
         A_t

      (horizontal, dynamical)

  ... -> E_t --> A_t --> E_{t+1} --> A_{t+1} --> ...
            ^                       |
            +----- impulses I_t ----+
\end{verbatim}

\end{minipage}
\end{center}

The \emph{vertical} arrows \( \Will \to E_t \to A_t \) tell us \emph{what} kind of entities we are dealing with and how they depend on each other ontologically:

\begin{itemize}[leftmargin=2em]
    \item \(E_t\) is always an element of \(\mathcal{E}(\Will)\): a manifestation of will.
    \item \(A_t\) is always an element of \(\mathcal{A}(E_t)\): an unfolding of will through that structure.
\end{itemize}

The \emph{horizontal} arrows, governed by \(F\) and \(G\), tell us \emph{how} these entities evolve and interact in time:

\begin{itemize}[leftmargin=2em]
    \item given \(E_t\) and \(I_t\), the system produces \(A_t\);
    \item given \(E_t\), \(A_t\), \(I_t\), and \(\Theta\), the system updates to \(E_{t+1}\);
    \item the new structure \(E_{t+1}\) then generates future actions \(A_{t+1}\), and so on.
\end{itemize}

The result is a closed feedback loop: actions reshape structure, which reshapes future actions, all under the continual influence of impulses from the environment.

\section{Preserving UTAM Asymmetry}

It is crucial that the introduction of horizontal dynamics does not erase or invert the asymmetry expressed by the UTAM axiom. At every time step \(t\), the triple \((\Will, E_t, A_t)\) still satisfies the vertical priority relation:
\[
    E_t \in \mathcal{E}(\Will), \qquad A_t \in \mathcal{A}(E_t).
\]

In other words:

\begin{itemize}[leftmargin=2em]
    \item \textbf{Every structure is still grounded in will.} The fact that \(E_t\) changes according to \(G\) does not make it independent of \(\Will\); it merely describes how one manifestation of will gives way to another over time.
    \item \textbf{Every action is still grounded in structure.} The fact that \(A_t\) depends on impulses \(I_t\) and on past actions does not break the dependence of action on structure; it enriches the way \(E_t\) is brought to bear in context.
\end{itemize}

Horizontal feedback thus operates entirely \emph{within} the space \(\mathcal{E}(\Will)\) of possible structures and the associated action spaces \(\mathcal{A}(E)\). It does not create or destroy \(\Will\); it traces \emph{trajectories} through the manifold of its manifestations.

This layered view---vertical ontology plus horizontal dynamics---is the backbone of the ONTO$\Sigma$ framework. In the chapters that follow, we will refine \(F\) and \(G\) into concrete algorithmic schemes (I$^2$C and $\Delta E$), introduce measures of coherence and drift, and reinterpret psychological, biological, and social phenomena as particular patterns in these coupled vertical--horizontal structures.


\chapter{Coherence and Drift as Dynamical Quantities}

With the vertical and horizontal layers of UTAM in place, we can now make precise two central dynamical notions that appeared informally in Part~I: \emph{coherence} and \emph{drift}. These are the key observables of ONTO$\Sigma$'s dynamic layer. Where the UTAM axiom
\[
    \Will \to \Struct \to \Act
\]
specifies ontological priority, coherence and drift quantify how well a particular manifestation \((E_t, A_t)\) is managing to express will under the impulses \(I_t\) it encounters.

\section{Defining Coherence \texorpdfstring{\(C_t\)}{Ct}}

At each time step \(t\), we would like to say, in a single symbol, how well the system is ``holding together'' in its current situation. We therefore introduce a coherence measure
\[
    C_t \in [0,1],
\]
interpreted as the degree of alignment between:

\begin{itemize}[leftmargin=2em]
    \item the current structure \(E_t\),
    \item the current pattern of impulses \(I_t\),
    \item and the current action \(A_t = F(E_t, I_t)\).
\end{itemize}

Intuitively:

\begin{itemize}[leftmargin=2em]
    \item \(C_t \approx 1\) means that the system's structure is well-tuned to its environment and its own behaviour; it is responding appropriately without undermining its own organisation.
    \item \(C_t \approx 0\) means that the system is badly misaligned; its structure no longer fits the demands placed upon it, and its actions are ineffective or self-destructive.
\end{itemize}

ONTO$\Sigma$ does not prescribe a single universal formula for \(C_t\); rather, it provides a \emph{role} that different concrete measures can play. Depending on context, \(C_t\) may be instantiated in several ways:

\paragraph{Information-theoretic interpretation.} In predictive processing and related frameworks, coherence can be tied to the quality of prediction. For example, one may let
\[
    C_t = f\bigl(-\text{prediction error}(E_t, I_t)\bigr)
\]
for some monotone mapping \(f\), or relate \(C_t\) to (the negative of) variational free energy. High coherence then corresponds to accurate, low-surprise modelling of incoming impulses by the current structure.

\paragraph{Control-theoretic interpretation.} In control and robotics, coherence can be related to stability and robustness. One might define \(C_t\) in terms of:

\begin{itemize}[leftmargin=2em]
    \item how close the system is to a desired manifold or invariant set,
    \item how robust the closed-loop behaviour is to perturbations,
    \item or how efficiently the system maintains key variables within viable bounds.
\end{itemize}

Here, high \(C_t\) means the controller--plant combination is functioning as intended under the current operating conditions.

\paragraph{Phenomenological interpretation.} From the first-person perspective, coherence often presents as a sense of ``rightness'' or fit: feeling attuned to one's situation, acting in ways that are consonant with one's deeper commitments, and not being internally at war. In such cases \(C_t\) is not directly measured, but \emph{inferred} from patterns of experience and behaviour. ONTO$\Sigma$ treats these lived qualities as legitimate indicators of how well \(\Will\)–\(\Struct_t\)–\(\Act_t\) are aligned.

Across these interpretations, the unifying idea is that \(C_t\) captures the \emph{current quality} of the system's directedness: how successfully will, in this particular structured form, is navigating the impulses it receives.

\section{Drift \texorpdfstring{\(D_t = 1 - C_t\)}{Dt = 1 - Ct}}

With coherence defined, \emph{drift} is introduced as its complement:
\[
    D_t = 1 - C_t.
\]

Where \(C_t\) measures alignment, \(D_t\) measures misalignment. A system with \(D_t \approx 0\) is, for the moment, well-matched to its environment and its own activity; a system with \(D_t \approx 1\) is deeply out of step.

It is useful to distinguish:

\begin{description}[leftmargin=2em]
    \item[Local drift.] Short-term deviations between structure, impulses, and actions. For example, a sudden change in inputs may temporarily lower \(C_t\) before adaptation restores it.
    \item[Long-term drift.] Persistent trends in which \(C_t\) remains low or declines over extended periods. This corresponds to chronic misalignment: a model trained on outdated data, an institution clinging to obsolete policies, a person living in a way that steadily erodes their integrity.
\end{description}

In many applications, the temporal profile of \(C_t\) and \(D_t\) is more informative than their instantaneous values. ONTO$\Sigma$ is especially concerned with regimes where, under changing conditions, \(D_t\) tends to increase unless effective structural adaptation is in place. This leads us to the Drift Law.

\section{The Drift Law as a Property of \texorpdfstring{\(G\)}{G}}

Recall the horizontal UTAM dynamics:
\begin{equation*}
\begin{aligned}
    A_t &= F(E_t, I_t), \\
    E_{t+1} &= G(E_t, A_t, I_t, \Theta).
\end{aligned}
\end{equation*}

The \textbf{Drift Law} is best understood as a qualitative property of the structure-updater \(G\) when the environment is non-stationary.

Suppose that:

\begin{enumerate}[leftmargin=2em]
    \item The environment is \textbf{non-stationary}: the distribution of impulses \(I_t\) changes over time. Informally, we write this as
    \[
        \frac{dI}{dt} \neq 0.
    \]
    \item The update rule \(G\) is \textbf{insensitive} to these changes, in the sense that the way it updates \(E_t\) does not significantly depend on the current impulses. Formally, we can express this as
    \[
        \frac{\partial G}{\partial I} \approx 0
    \]
    or, in a more extreme case, as adaptation parameters \(\Theta\) being effectively frozen.
\end{enumerate}

Under these conditions, we should expect that the structure \(E_t\) will gradually become less well-matched to the actual impulses it encounters. Since \(E_t\) is not updating in a way that tracks the changing \(I_t\), the alignment between \(E_t\), \(I_t\), and the resulting \(A_t\) will degrade.

In terms of coherence and drift, the Drift Law can then be stated qualitatively as:
\[
    \frac{dC_t}{dt} \leq 0, \qquad
    \frac{dD_t}{dt} \geq 0,
\]
for as long as the environment continues to change and \(G\) remains insensitive.

This is not a theorem about a specific equation; it is a \emph{structural insight}: in any non-stationary world, a system whose internal structure does not adequately register and integrate the changes in its impulses will \emph{inevitably} drift out of coherence. The precise rate and pattern of this drift depend on the details of \(F\), \(G\), and \(C_t\), but the direction of change is robust.

When we later introduce explicit $\Delta E$ dynamics, we will see that these can be understood as attempts to make \(\partial G / \partial I\) non-negligible in a way that tends to increase coherence rather than decrease it.

\section{Drift Across Levels}

One of the strengths of ONTO$\Sigma$ is that coherence and drift are defined at a level of abstraction that applies across many kinds of systems. Here we briefly indicate how drift appears at different levels, foreshadowing the domain-specific discussions in Part~III.

\paragraph{Toy models.} In simple simulation models, we might let the environment be characterized by a slowly drifting parameter \(k_t\), while the system maintains an internal estimate \(E_t\) of this parameter. If \(E_t\) is held fixed while \(k_t\) drifts, coherence (e.g.\ how well \(E_t\) predicts outcomes) will decline: a direct visualization of the Drift Law. If \(E_t\) is updated via a $\Delta E$ rule tuned to track \(k_t\), coherence can be maintained.

\paragraph{Learning systems.} In machine learning, models are typically trained on a particular data distribution and then deployed in environments whose statistics may change. When the training distribution and deployment distribution diverge (\(dI/dt \neq 0\)) and no online adaptation is used (\(\partial G/\partial I \approx 0\)), performance degrades. Here, \(C_t\) might be instantiated as accuracy or negative loss; drift corresponds to the familiar phenomenon of performance decay under distribution shift.

\paragraph{Human cognition and behaviour.} In human life, drift shows up when a person's internal structures---their habits, beliefs, narratives---fail to update in the face of new realities. The career that once made sense becomes hollow; the coping strategies that once worked begin to fail; actions that previously expressed a coherent self start to feel mechanical or self-betraying. Coherence here encompasses both objective fit (ability to navigate the world) and subjective sense of meaning and authenticity.

\paragraph{Institutions and societies.} At larger scales, legal systems, economies, and cultural norms drift when they lag behind technological, ecological, or demographic changes. An institution that does not revise its structures (laws, procedures, incentives) in response to new impulses will experience growing incoherence: policies that no longer achieve their intended effects, rules that are widely flouted, narratives that no longer resonate.

In all of these cases, the same pattern recurs: non-stationary impulses, insufficient structural adaptation, declining coherence. The details differ, but the underlying structure is the same. This is precisely what UTAM aims to capture: a single, unified language for directedness, coherence, and drift across levels. The next chapters will deepen this picture by specifying concrete adaptation schemes (I$^2$C and $\Delta E$) and by introducing the geometry of the volitional phase space in which these dynamics unfold.


\chapter{Algorithmic Schemes: I$^2$C and \texorpdfstring{$\Delta E$}{ΔE}}

The horizontal UTAM equations
\[
\begin{aligned}
    A_t &= F(E_t, I_t), \\
    E_{t+1} &= G(E_t, A_t, I_t, \Theta),
\end{aligned}
\]
describe the temporal evolution of actions and structures in abstract form. To make them operational, we now introduce two algorithmic schemes that instantiate \(F\) and \(G\) in a wide range of adaptive systems:

\begin{itemize}[leftmargin=2em]
    \item the \emph{I$^2$C cycle} (Impulse $\to$ Interpretation $\to$ Coherence), which implements the action-generator \(F\);
    \item the \(\Delta E\) principle, which implements the structure-updater \(G\) as a coherence-preserving adaptation rule.
\end{itemize}

These schemes are not arbitrary designs; they are minimal patterns that any system must approximate if it is to maintain coherence under drift.

\section{The I\texorpdfstring{$^2$}{²}C Cycle (Impulse $\to$ Interpretation $\to$ Coherence)}

At each time step, an adaptive system faces three basic tasks:

\begin{enumerate}[leftmargin=2em]
    \item receive an \emph{impulse} from its environment;
    \item \emph{interpret} this impulse in light of its current structure;
    \item choose an \emph{action} that maintains or restores coherence.
\end{enumerate}

We model this as the I$^2$C cycle:

\paragraph{Impulse.} The environment generates impulses \(I_t\) according to some (possibly non-stationary) distribution:
\[
    I_t \sim P_{\mathrm{env}}(\cdot).
\]

\paragraph{Interpretation.} The system transforms the raw impulse into an internal representation \(Z_t\) using its current structure \(E_t\):
\[
    Z_t = \mathrm{Enc}(I_t, E_t),
\]
where \(\mathrm{Enc}\) is an \emph{encoder} that may be realized as a neural network, a Bayesian filter, a symbolic parser, or any other mapping from input and structure to a meaningful internal state.

\paragraph{Coherence/action.} Given the interpretation \(Z_t\) and the structure \(E_t\), the system chooses an action \(A_t\) by maximizing a coherence functional \(C\):
\[
    A_t = \arg\max_{A \in \mathcal{A}(E_t)} C(Z_t, A, E_t).
\]
Here \(C(Z_t, A, E_t)\) quantifies how well a candidate action \(A\) fits the interpreted impulse under the current structure. The system selects the action that yields the highest coherence.

Putting these steps together, we obtain an explicit form for the action-generator:
\[
    F(E_t, I_t) = \arg\max_{A \in \mathcal{A}(E_t)} 
        C\bigl(\mathrm{Enc}(I_t, E_t), A, E_t\bigr).
\]

The I$^2$C cycle thus realizes \(F\) as a two-stage process: encoding the impulse, then choosing an action that best maintains coherence. It is ``minimal'' in the sense that any system that \emph{does not} in some way interpret its inputs and evaluate the coherence of candidate actions will, under drift, be unable to maintain alignment between \(\Will\), \(E_t\), and \(A_t\).

\section{\texorpdfstring{$\Delta E$}{ΔE}: Structure Update to Preserve Coherence}

The I$^2$C cycle, by itself, can select coherent actions \emph{given} a structure \(E_t\). But as we have seen, in a non-stationary environment the structure itself must adapt if coherence is to be preserved. This is the role of the \(\Delta E\) principle: to update \(E_t\) in a way that tends to increase coherence over time.

\subsection*{The $\Delta E$ principle}

The core idea of $\Delta E$ is:

\begin{quote}
The purpose of structural update is not to optimize a fixed external objective, but to \emph{preserve or improve coherence} between the system and its world under changing conditions.
\end{quote}

Formally, we posit an update of the form
\[
    E_{t+1} = G(E_t, A_t, I_t, \Theta),
\]
where \(G\) is constructed so that, in expectation, the coherence \(C_{t+1}\) is not lower than \(C_t\), or at least decreases as slowly as possible given constraints.

A canonical way to implement this is via a gradient-based rule:

\begin{equation}
    E_{t+1} = E_t + \eta \, \nabla_{E} C(E_t, I_t, A_t),
    \label{eq:delta-e-gradient}
\end{equation}
where:

\begin{itemize}[leftmargin=2em]
    \item \(C(E_t, I_t, A_t)\) is the coherence functional evaluated at the actual action,
    \item \(\nabla_E C\) is its gradient with respect to the structure \(E_t\),
    \item \(\eta > 0\) is a learning rate or adaptation speed.
\end{itemize}

Equation~\eqref{eq:delta-e-gradient} performs \emph{coherence gradient ascent}: it adjusts the structure in the direction that locally increases coherence, given the recent impulse and action.

Crucially, the \(\Delta E\) principle is \emph{broader} than this specific gradient form. It refers to any update rule that systematically increases coherence \emph{in expectation}. Concrete realizations include:

\begin{itemize}[leftmargin=2em]
    \item \textbf{Gradient-based learning.} As in \eqref{eq:delta-e-gradient}, where \(\nabla_E C\) is computed analytically or approximated (e.g.\ through stochastic gradients).
    \item \textbf{Hebbian / synaptic rules.} ``Cells that fire together wire together'' can be seen as a local approximation to coherence maximization: structures that consistently support successful responses are strengthened.
    \item \textbf{Bayesian model updating.} Updating a generative model \(E_t\) by Bayes' rule in light of new data \(I_t\) tends, under mild conditions, to improve predictive coherence between model and environment.
\end{itemize}

In all these cases, the essence of \(\Delta E\) is the same: the system changes \emph{how it is organized} in order to remain a viable manifestation of \Will{} in the face of drift.

\section{Relation to Existing Theories}

The I$^2$C and $\Delta E$ schemes are intentionally generic. They are meant to capture the structural commonality behind a range of existing theories and algorithms.

\subsection*{Predictive processing and the Free Energy Principle}

In predictive processing frameworks, the brain is modelled as maintaining a generative model \(E_t\) that predicts incoming sensory data. Perception involves adjusting internal states to minimize prediction error; action involves changing the world to make predictions come true; learning involves updating the generative model itself.

In this language:

\begin{itemize}[leftmargin=2em]
    \item The encoding \(Z_t = \mathrm{Enc}(I_t, E_t)\) corresponds to perceptual inference: estimating latent causes given data and a model.
    \item The coherence functional \(C\) is related to (the negative of) free energy or prediction error: higher \(C\) means better prediction/lower surprise.
    \item The choice of action \(A_t\) to reduce expected free energy is an instance of the I$^2$C ``coherence/action'' step.
    \item Learning the model (updating \(E_t\)) to improve predictions is a form of $\Delta E$.
\end{itemize}

ONTO$\Sigma$ does not compete with these ideas; it provides an ontological and dynamical umbrella under which they can be seen as one realization of I$^2$C + $\Delta E$.

\subsection*{Reinforcement learning: reward vs.\ coherence}

In reinforcement learning (RL), an agent selects actions to maximize expected cumulative reward. Policies and value functions are updated to better approximate the optimal mapping from states to actions.

In UTAM terms:

\begin{itemize}[leftmargin=2em]
    \item Reward can be seen as one possible ingredient of the coherence functional \(C\), but not the only one. A system could sacrifice long-term structural integrity in order to maximize short-term reward; this would be incoherent in the ONTO$\Sigma$ sense even if it were ``optimal'' in the narrow RL sense.
    \item Policy updates in RL (gradient ascent on expected reward) resemble $\Delta E$ updates, but with a fixed external objective. Genuine $\Delta E$ updates, in contrast, may allow \(C\) itself to change as the system's structure and context evolve.
\end{itemize}

This suggests a shift of emphasis: from maximizing externally prescribed rewards to maintaining and enhancing internal coherence as the primary imperative of directed systems.

\subsection*{Adaptive and robust control}

In control theory, adaptive and robust controllers adjust their parameters in response to changes in plant dynamics or disturbances, seeking to maintain stability and performance.

Here:

\begin{itemize}[leftmargin=2em]
    \item The controller's parameter vector is part of \(E_t\).
    \item The plant output and reference signals play the role of impulses \(I_t\).
    \item The control signal is the action \(A_t\).
    \item Adaptive update laws for the controller parameters are instances of $\Delta E$, often designed to maintain stability margins or minimize tracking error, which can be viewed as concrete forms of coherence.
\end{itemize}

Classical adaptive control thus fits neatly into the I$^2$C + $\Delta E$ framework, albeit often with a narrower notion of coherence focused on tracking a setpoint.

\section{Toy Example: A Robot on Drifting Terrain}

To make these ideas more tangible, consider a simple mobile robot tasked with walking across a terrain whose properties change slowly over time.

\subsection*{Setup}

\begin{itemize}[leftmargin=2em]
    \item The robot's structure \(E_t\) encodes its \emph{gait model}: parameters determining how joint angles and forces are coordinated to produce locomotion.
    \item The impulses \(I_t\) are sensor readings: joint positions, contact forces, inertial measurements, perhaps visual cues about the surface ahead.
    \item The actions \(A_t\) are motor commands: torques or position targets sent to actuators.
\end{itemize}

Initially, the robot is trained to walk on flat, high-friction terrain. Over time, the environment drifts: the ground becomes softer, more uneven, or more slippery.

\subsection*{I$^2$C in the robot}

At each time step:

\begin{itemize}[leftmargin=2em]
    \item \textbf{Impulse.} The robot receives \(I_t\) from its sensors, reflecting the current terrain and its own body state.
    \item \textbf{Interpretation.} An encoder \(\mathrm{Enc}(I_t, E_t)\) transforms these raw signals into an internal state \(Z_t\), perhaps estimating foot contact phases, slip, or balance margins.
    \item \textbf{Coherence/action.} The robot selects motor commands
    \[
        A_t = \arg\max_{A} C(Z_t, A, E_t),
    \]
    where \(C\) measures, for example, expected stability and forward progress without exceeding actuator limits.
\end{itemize}

This is an instance of the I$^2$C cycle: the action chosen is the one that, given the current interpretation and gait model, is most likely to preserve coherent locomotion.

\subsection*{$\Delta E$ vs.\ fixed controller}

Now compare two scenarios:

\begin{description}[leftmargin=2em]
    \item[Fixed controller (no $\Delta E$).] The robot's gait parameters \(E_t\) are held fixed at their initial values. As the terrain drifts away from the training conditions, the same I$^2$C pipeline produces actions that are increasingly inappropriate. Coherence \(C_t\) declines: the robot slips, stumbles, or falls more often. This is a direct illustration of the Drift Law.
    \item[Adaptive controller ($\Delta E$ active).] The robot updates its gait parameters according to a $\Delta E$ rule, for example
    \[
        E_{t+1} = E_t + \eta \, \nabla_E C(E_t, I_t, A_t),
    \]
    where \(C\) now reflects not only immediate stability but also energy efficiency and long-term viability. As the terrain characteristics drift, the robot's gait gradually changes to maintain a high level of coherence. Performance degrades more slowly or not at all.
\end{description}

This toy example makes concrete the difference between:

\begin{itemize}[leftmargin=2em]
    \item systems that treat their structure as fixed and optimize within that structure (susceptible to drift), and
    \item systems that treat structural adaptation as a first-class process aimed at preserving coherence (instances of I$^2$C + $\Delta E$).
\end{itemize}

In later sections, we will see how similar patterns play out in learning systems, human cognition, and social institutions. The point here is that the same abstract schemes---I$^2$C for action generation and $\Delta E$ for structural update---provide a compact, UTAM-consistent way to describe adaptive behaviour across domains.

\chapter{Volitional Phase Space: Geometric Interpretation}

The dynamical equations of ONTO$\Sigma$ describe how structures \(E_t\) and actions \(A_t\) evolve over time under impulses \(I_t\). To gain a more intuitive grasp of these dynamics, it is useful to recast them in geometric terms. In this chapter we introduce the notion of a \emph{volitional phase space}: a space of possible structures and impulses equipped with a \emph{coherence landscape} on which the system moves.

\section{Structural Manifold and Volitional Phase Space}

We begin by formalizing the space of structures.

\subsection*{Structural manifold}

Let
\[
    \mathcal{M} := \mathcal{E}(\Will)
\]
denote the \textbf{structural manifold}: the space of all structures that can manifest will. Each point \(E \in \mathcal{M}\) corresponds to a particular organisation of directedness: a set of physical laws, a neural architecture, a policy network, a pattern of norms, and so on.

We do not assume a specific dimensionality or topology for \(\mathcal{M}\), only that it can be treated as a (possibly high-dimensional) manifold on which local notions such as gradients make sense.

\subsection*{Volitional phase space}

Let \(\mathcal{I}\) denote the space of possible impulses, i.e.\ the types of environmental input a system can receive. Then we define the \textbf{volitional phase space} as
\[
    \mathcal{V} := \mathcal{M} \times \mathcal{I}.
\]

A point \((E, I) \in \mathcal{V}\) specifies both:

\begin{itemize}[leftmargin=2em]
    \item a particular structure \(E\) (a way in which \Will{} is organized), and
    \item a particular impulse configuration \(I\) (what the environment is ``doing'' to the system).
\end{itemize}

The horizontal UTAM dynamics can then be thought of as generating trajectories in \(\mathcal{V}\), with the structural component evolving via $\Delta E$ and the impulse component evolving according to the environment.

\section{Coherence Landscape}

On this volitional phase space, we introduce a scalar field that captures how well each structure--impulse combination fits together.

\subsection*{Coherence as a scalar field}

Define the \textbf{coherence landscape} as a function
\[
    C : \mathcal{M} \times \mathcal{I} \to [0,1], \qquad
    (E, I) \mapsto C(E, I),
\]
where \(C(E, I)\) measures the coherence that would result if the system with structure \(E\) encountered impulse \(I\) and acted according to its I$^2$C dynamics.

Intuitively:

\begin{itemize}[leftmargin=2em]
    \item Points \((E, I)\) where \(C(E, I)\) is close to \(1\) correspond to situations where the structure \(E\) is well-matched to the impulse \(I\); actions generated from this pair will maintain or enhance the system's integrity.
    \item Points where \(C(E, I)\) is close to \(0\) correspond to situations of deep misfit; actions generated from this pair will tend to be ineffective or self-undermining.
\end{itemize}

We can visualize (in low dimensions) \(C(E, I)\) as a landscape:

\begin{itemize}[leftmargin=2em]
    \item \textbf{Peaks} (local maxima) represent high-coherence configurations: attractor-like combinations of structure and impulse where the system functions smoothly.
    \item \textbf{Valleys} (local minima) represent low-coherence configurations: regions of confusion, instability, or breakdown.
\end{itemize}

The coherence landscape is not static in general; as the environment and the system's own internal criteria shift, the shape of \(C\) may change. For the moment, however, we treat \(C\) as a fixed function on \(\mathcal{V}\) in order to describe the basic geometry of $\Delta E$ dynamics.

\section{Trajectories and Gradient Flow}

Given the coherence landscape, the $\Delta E$ principle can be interpreted as a gradient flow on the structural manifold.

\subsection*{Continuous-time approximation}

In continuous time, a natural idealization of $\Delta E$ is
\begin{equation}
    \frac{dE}{dt} = \eta \, \nabla_E C(E, I) + \xi(t),
    \label{eq:gradient-flow}
\end{equation}
where:

\begin{itemize}[leftmargin=2em]
    \item \(\nabla_E C(E, I)\) is the gradient of coherence with respect to the structural coordinates,
    \item \(\eta > 0\) is an adaptation rate or learning rate,
    \item \(\xi(t)\) is a noise term capturing stochasticity, modelling error, or unobserved influences.
\end{itemize}

If we temporarily ignore noise and treat \(I\) as fixed, \eqref{eq:gradient-flow} describes a trajectory on \(\mathcal{M}\) that tends to move uphill on the coherence landscape: the system modifies its structure in the direction that increases \(C(E, I)\).

\subsection*{Role of learning rate and noise}

The parameters \(\eta\) and \(\xi(t)\) have clear geometric interpretations:

\begin{itemize}[leftmargin=2em]
    \item The \textbf{learning rate} \(\eta\) controls the step size of the flow. Small \(\eta\) leads to slow, cautious adaptation that closely tracks the gradient; large \(\eta\) leads to rapid but potentially unstable jumps that may overshoot coherent regions.
    \item The \textbf{noise} term \(\xi(t)\) can either hinder or help. Too much noise may prevent the system from ascending the coherence gradient; a moderate amount can help escape shallow local maxima and explore the landscape more broadly.
\end{itemize}

From this perspective, the temporal evolution of \(E_t\) is a path on the structural manifold \(\mathcal{M}\) shaped by both the geometry of the coherence landscape and the adaptation dynamics encoded by $\Delta E$.

\section{Drift as Landscape Deformation}

So far we have treated the impulse \(I\) as fixed when describing gradient flow. In reality, however, the environment is typically non-stationary: the impulse component changes over time according to some dynamics
\[
    \frac{dI}{dt} = f(I, \text{environment}).
\]

\subsection*{Changing impulses, changing landscape}

Because \(C\) depends on both \(E\) and \(I\), changes in \(I\) effectively \emph{deform} the coherence landscape:
\[
    C_t(E) := C(E, I_t).
\]

As \(I_t\) evolves, the peaks and valleys of \(C_t(E)\) shift:

\begin{itemize}[leftmargin=2em]
    \item Coherent regions (high \(C_t\)) may move in \(\mathcal{M}\) or change shape.
    \item Previously coherent structures \(E\) may find themselves in lower-coherence regions as the environment drifts.
\end{itemize}

\subsection*{Standing still as sliding downhill}

If the system's structure fails to adapt---for example, if \(dE/dt \approx 0\) because \(G\) is insensitive to \(I\) or adaptation parameters \(\Theta\) are frozen---then the point \((E_t, I_t)\) moves in \(\mathcal{V}\) purely because of changes in \(I_t\). Geometrically, this corresponds to \emph{sliding across a changing landscape} without adjusting one's position.

In terms of coherence, this is precisely the situation described by the Drift Law:

\begin{itemize}[leftmargin=2em]
    \item If \(\frac{dI}{dt} \neq 0\) (the environment changes),
    \item and \(\frac{dE}{dt} \approx 0\) or \(\frac{\partial G}{\partial I} \approx 0\) (the structure does not respond),
    \item then, generically, \(C_t(E_t) = C(E_t, I_t)\) will tend to decrease over time:
    \[
        \frac{dC_t}{dt} \leq 0.
    \]
\end{itemize}

In other words, ``standing still'' in \(\mathcal{M}\) while the impulse component changes is equivalent, from the system's perspective, to \emph{sliding downhill} on the coherence landscape. Drift is thus naturally interpreted as the combination of environmental deformation of \(C\) and insufficient structural motion along its gradients.

In the next sections (not included here), we will deepen this geometric picture by connecting specific psychological states to features of the coherence landscape and by examining how meta-structural processes can reshape the landscape itself, expanding the accessible regions of \(\mathcal{E}(\Will)\).


\section{Psychological States as Attractors}

The coherence landscape on the volitional phase space \(\mathcal{V} = \mathcal{M} \times \mathcal{I}\) provides a natural way to interpret familiar psychological states as \emph{attractor structures} in which the system tends to settle (or become stuck). Here we sketch four such regimes.

\subsection*{Flow: deep, wide coherence basins}

In \emph{flow} states, experience is often described as effortless, focused, and intrinsically rewarding. In ONTO$\Sigma$ terms, this corresponds to the system occupying a deep, wide basin of high coherence:

\begin{itemize}[leftmargin=2em]
    \item The coherence landscape \(C(E, I)\) has a broad region where \(C\) is close to 1.
    \item The gradient \(\nabla_E C\) within this basin is gentle: small deviations in structure are naturally corrected by the dynamics without violent oscillations.
    \item The environment's impulses \(I_t\) vary within a range that the current structure \(E_t\) can handle without being pushed out of the basin.
\end{itemize}

Flow is thus interpreted as the system temporarily inhabiting a robust attractor in \(\mathcal{M}\), where both local perturbations and moderate drift in \(I_t\) are absorbed without large drops in coherence.

\subsection*{Burnout: low, flat basins and weak gradients}

\emph{Burnout} and related states (exhaustion, depersonalization) correspond to a different attractor morphology:

\begin{itemize}[leftmargin=2em]
    \item The system has settled into a region where \(C(E, I)\) is relatively low but locally \emph{flat}: \(|\nabla_E C|\) is small.
    \item Because the gradient is weak, $\Delta E$ updates are ineffective; even if the system ``tries'' to adapt, the structure does not move appreciably toward higher-coherence regions.
    \item Subjectively, this appears as a sense of stagnation, futility, or being stuck: actions no longer feel connected to meaningful improvement.
\end{itemize}

In this picture, burnout is not merely a lack of energy but a structural regime in which the coherence landscape offers few usable gradients for adaptation.

\subsection*{Anxiety: rugged landscape and rapid changes}

\emph{Anxiety} can be interpreted as the experience of navigating a rugged, rapidly changing coherence landscape:

\begin{itemize}[leftmargin=2em]
    \item \(C(E, I)\) has many narrow peaks and deep valleys: small changes in \(E\) or \(I\) produce large swings in coherence.
    \item The environment-driven dynamics \(dI/dt\) are fast relative to the system's adaptation rate \(\eta\), causing the coherent regions to move quickly.
    \item The system thus finds itself repeatedly climbing steep local gradients only to have the landscape shift, leading to frequent drops in \(C_t\).
\end{itemize}

Phenomenologically, this corresponds to hypervigilance, instability, and a sense that no stable, coherent way of being can be maintained for long.

\subsection*{Aging: slow deformation of attractor structure}

\emph{Aging}---both biological and psychological---can be viewed as a slow deformation of the coherence landscape and of the system's typical trajectory within it:

\begin{itemize}[leftmargin=2em]
    \item Over long timescales, the shape of \(C(E, I)\) changes due to cumulative internal changes (wear, learning history, path dependencies) and external drifts.
    \item Attractors that were once deep and accessible may become shallower or shift location; new attractors may appear that reflect different capacities or constraints.
    \item The system's $\Delta E$ dynamics may slow down (effective \(\eta\) decreases), making it harder to reach newly emergent high-coherence regions.
\end{itemize}

From this perspective, aging is neither purely decline nor purely growth; it is a reconfiguration of the coherence landscape and of the structural pathways by which the system can move through it.

\section{Meta-Structural Dynamics}

So far we have treated $\Delta E$ as adjustments of structure \emph{within} a given coherence landscape. However, some processes appear to change the \emph{landscape itself}: they alter what kinds of structures are possible or salient. ONTO$\Sigma$ refers to these as \emph{meta-structural dynamics}.

\subsection*{Meta-will and landscape reshaping}

Meta-structural processes are those by which a system:

\begin{itemize}[leftmargin=2em]
    \item modifies its own adaptation rules (changing \(G\) or \(\eta\)),
    \item introduces new structural dimensions or constraints,
    \item or reorganizes the criteria that define coherence \(C\).
\end{itemize}

Examples include:

\begin{description}[leftmargin=2em]
    \item[Conscious reflection.] A person can step back from their current habits and narratives, \emph{think about} how they adapt, and deliberately choose new strategies. This is not just moving within \(\mathcal{M}\), but altering the effective geometry of the local coherence landscape.
    \item[Cultural evolution.] Shared languages, norms, and institutions create new structural possibilities and coherence criteria that did not exist at the individual level. Entire regions of \(\mathcal{E}(\Will)\) become newly accessible or newly salient.
    \item[Technological innovation.] Tools and infrastructures (from writing to AI) expand the range of structures that can be instantiated and maintained, effectively opening new dimensions in \(\mathcal{M}\) and changing how \(C(E, I)\) behaves.
\end{description}

In UTAM terms, these processes correspond to an \emph{expansion of the accessible region} of \(\mathcal{E}(\Will)\) without altering \Will{} itself. They are higher-order manifestations of will: will shaping not just actions and structures, but the very space of possible structures.

\section{Mathematical Sketch}

We can summarize the geometric and dynamical picture developed so far in a compact mathematical form.

\subsection*{Coupled dynamics on \(\mathcal{V}\)}

Let \(\mathcal{M} = \mathcal{E}(\Will)\) be the structural manifold, \(\mathcal{I}\) the impulse space, and \(\mathcal{V} = \mathcal{M} \times \mathcal{I}\) the volitional phase space. The dynamics are given by:

\begin{align}
    \frac{dE}{dt} &= \eta \, \nabla_E C(E, I) + \xi(t),
    \label{eq:coupled-dE}\\[0.5em]
    \frac{dI}{dt} &= f(I, \text{environment}),
    \label{eq:coupled-dI}
\end{align}
where:

\begin{itemize}[leftmargin=2em]
    \item \(C(E, I) : \mathcal{M} \times \mathcal{I} \to [0,1]\) is the coherence landscape,
    \item \(\eta\) is an adaptation rate,
    \item \(\xi(t)\) is a noise term,
    \item \(f\) encodes the environment's dynamics (which may or may not depend on the system itself).
\end{itemize}

Together, \eqref{eq:coupled-dE} and \eqref{eq:coupled-dI} define a \emph{time-dependent vector field} on \(\mathcal{V}\). Attractors in this field---sets toward which trajectories converge---correspond to coherence basins in which the system tends to remain. As \(f\) changes (e.g.\ under slow environmental drift), these attractors can move, merge, or disappear.

Meta-structural dynamics would appear, in this sketch, as changes to:

\begin{itemize}[leftmargin=2em]
    \item the definition of \(C\) (reshaping the landscape),
    \item the structure of \(\mathcal{M}\) itself (adding or removing dimensions),
    \item or the parameters of the $\Delta E$ flow (changing \(\eta\), the form of \(\nabla_E C\), or the statistics of \(\xi\)).
\end{itemize}

ONTO$\Sigma$ does not commit to a specific analytic form for \(C\) or \(f\); the point of this sketch is to show how the concept of a volitional phase space can be made mathematically concrete.

\section{Example: The Wim Hof Practitioner}

To illustrate these ideas in a more experiential context, consider a person training in the Wim Hof Method (WHM), which combines breathing exercises and cold exposure.

\subsection*{Early incoherent responses}

At the outset:

\begin{itemize}[leftmargin=2em]
    \item The person's physiological and psychological structure \(E_0\) includes standard autonomic responses to cold: vasoconstriction, shivering, heightened stress responses.
    \item The impulse \(I_t\) corresponding to cold water immersion initially places the system in a region \((E_0, I_{\mathrm{cold}})\) where \(C(E_0, I_{\mathrm{cold}})\) is low: the experience is overwhelming, actions are disorganized (panic, hyperventilation), and coherence is temporarily disrupted.
\end{itemize}

In the coherence landscape, \((E_0, I_{\mathrm{cold}})\) lies in or near a low-coherence valley.

\subsection*{$\Delta E$ practice and movement into a high-coherence basin}

Through repeated, structured practice:

\begin{itemize}[leftmargin=2em]
    \item The practitioner engages in deliberate breathing patterns, graded exposure, and cognitive reframing. These act as a combination of impulses \(I_t\) and internally generated actions \(A_t\) that drive $\Delta E$ updates.
    \item Over time, the autonomic and cognitive structure \(E_t\) changes: breathing patterns become more efficient, vascular responses adapt, fear responses are modulated.
    \item In the volitional phase space, the system moves from \((E_0, I_{\mathrm{cold}})\) toward a new region \((E^\ast, I_{\mathrm{cold}})\) where \(C(E^\ast, I_{\mathrm{cold}})\) is high: the same cold impulse now elicits calm, controlled breathing and a sense of clarity.
\end{itemize}

In geometric terms, $\Delta E$ has carried the system into a new coherence basin associated with cold exposure.

\subsection*{Meta-structural mastery}

With continued practice and reflection, the practitioner may reach a stage where:

\begin{itemize}[leftmargin=2em]
    \item They can \emph{intentionally access} this high-coherence basin: by adjusting breathing and attention, they can move their structure \(E_t\) into a configuration appropriate for cold or stress even before the impulse arrives.
    \item They can \emph{teach} the method to others, articulating the underlying patterns and principles. This contributes to cultural evolution: new structural possibilities (ways of relating to cold, stress, and breath) become part of the shared coherence landscape of a community.
\end{itemize}

This transition from individual adaptation to transmissible know-how is an instance of meta-structural dynamics: the system is no longer merely flowing along a fixed landscape, but helping to reshape the landscape itself for others. The accessible region of \(\mathcal{E}(\Will)\) associated with human physiological and psychological responses has been expanded, without any change to \Will{} as such.

In this way, the Wim Hof example concretely illustrates the full arc of ONTO$\Sigma$'s geometric vocabulary: low-coherence valleys, gradient-driven $\Delta E$ ascent, stable attractors (flow-like states), and meta-structural processes that propagate and transform these attractors at larger scales.

% ================================================
% Part III – Applications and Empirical Directions
% ================================================

\part{Applications and Empirical Directions}

\chapter{Neuroscience and Physiology}

The ONTO$\Sigma$ framework is not intended to remain at the level of abstract metaphysics and general dynamics. One of its central claims is that coherence \(C_t\) and drift \(D_t\) correspond to empirically measurable patterns in concrete systems. In this chapter we outline how UTAM can be brought into contact with contemporary neuroscience and physiology, focusing on neural coherence, maladaptive drift, and candidate experimental paradigms.

\section{Neural Coherence as \texorpdfstring{$C_t$}{C\_t}}

At the neural level, ONTO$\Sigma$ suggests that the coherence \(C_t\) of a cognitive--physiological system should be reflected in patterns of \emph{neural coordination}. While the exact mapping is open to empirical refinement, several candidate measures stand out:

\subsection*{Phase synchrony and oscillatory coupling}

A large body of work characterizes communication between brain regions in terms of oscillatory activity and phase synchrony. In UTAM terms, we can interpret:

\begin{itemize}[leftmargin=2em]
    \item \textbf{Local coherence} as consistent phase relationships within and between task-relevant regions during a given cognitive or affective state.
    \item \textbf{Global coherence} as patterns of long-range synchronization that unify distributed processes into a functional whole (e.g.\ fronto-parietal coupling during attention).
\end{itemize}

High \(C_t\) would then correspond, in many contexts, to organised synchrony patterns that support efficient information routing, while low \(C_t\) would be associated with either excessive desynchronization (fragmentation) or rigid, maladaptive coupling (overbinding).

\subsection*{Functional connectivity and network structure}

Beyond oscillations, functional connectivity analyses (e.g.\ correlation, Granger causality, directed information measures) provide a network-level view of how brain regions co-activate and influence each other. In ONTO$\Sigma$:

\begin{itemize}[leftmargin=2em]
    \item The evolving connectivity graph is part of the structure \(E_t\): a dynamic instantiation of how will is organized at the neural level.
    \item Coherence \(C_t\) can be tied to connectivity patterns that support stable task performance, flexible reconfiguration, and resilience to perturbation.
\end{itemize}

For example, a brain in a high-coherence state might show:

\begin{itemize}[leftmargin=2em]
    \item modular yet integrated network structure,
    \item efficient small-world properties,
    \item and context-appropriate engagement of control hubs.
\end{itemize}

A low-coherence state, by contrast, might show either overly fragmented networks (loss of integration) or hyper-connected patterns that prevent flexible reconfiguration.

\subsection*{Physiological coherence}

Coherence in ONTO$\Sigma$ is not purely neural. It also encompasses \emph{autonomic and physiological} variables:

\begin{itemize}[leftmargin=2em]
    \item heart-rate variability (HRV) as a marker of flexible autonomic regulation,
    \item respiratory patterns and their coupling to neural rhythms,
    \item endocrine signals (e.g.\ cortisol) reflecting stress-load and adaptation.
\end{itemize}

A plausible working hypothesis is that high \(C_t\) corresponds to coordinated patterns across these levels: neural networks, ANS, and bodily systems functioning in a mutually supportive regime.

\section{Drift and Maladaptation in the Brain}

If coherence \(C_t\) has neural and physiological signatures, then drift \(D_t = 1 - C_t\) should correspond to \emph{maladaptive regimes} where structure \(E_t\) fails to keep up with environmental and internal demands. Several conditions can be reinterpreted in this light.

\subsection*{Chronic stress}

Under persistent stressors, the environment's impulse stream \(I_t\) is effectively non-stationary and biased toward threat. If neural and physiological structures do not adapt in a way that restores coherence---for example, if coping strategies remain rigid, social context does not change, or recovery periods are insufficient---then:

\begin{itemize}[leftmargin=2em]
    \item coherence \(C_t\) gradually decreases,
    \item baseline arousal increases,
    \item and stress-related circuitry (e.g.\ amygdala, HPA axis) becomes chronically over-engaged.
\end{itemize}

This is a paradigmatic case of long-term drift: the system's default structure \(E_t\) no longer fits its actual impulse statistics, yet $\Delta E$ processes are either blocked or misdirected.

\subsection*{Burnout}

Burnout can be seen as a \emph{neurocognitive attractor} in which the coherence landscape has become shallow and flat around a low-\(C_t\) configuration:

\begin{itemize}[leftmargin=2em]
    \item stable connectivity patterns support repetitive, obligation-driven behaviour,
    \item reward and salience systems are blunted,
    \item and $\Delta E$ dynamics (learning, reframing, exploration) are weakened.
\end{itemize}

From the ONTO$\Sigma$ perspective, burnout is not just ``too much stress'', but a chronic drift regime in which the brain's structure has adapted \emph{toward} a low-coherence basin that is hard to escape with ordinary updates.

\subsection*{Post-traumatic stress disorder (PTSD)}

In PTSD, extreme impulses \(I_t\) (trauma) can induce abrupt and profound changes in structure \(E_t\):

\begin{itemize}[leftmargin=2em]
    \item threat-detection circuits become hypersensitized,
    \item contextual modulation (e.g.\ hippocampal input) is weakened,
    \item and default networks may be reorganised around hypervigilance and intrusive recall.
\end{itemize}

Coherence becomes highly context-dependent: in genuinely dangerous situations, the new structure may be locally coherent; in safe contexts, it produces inappropriate hyperarousal and intrusive memories. This is a form of \emph{maladapted drift}: the system has reorganized, but into a structural regime that is globally incoherent with its actual environment.

In all these cases, ONTO$\Sigma$ suggests that:

\begin{itemize}[leftmargin=2em]
    \item \(C_t\) and \(D_t\) can be operationalized via neural and physiological markers,
    \item and interventions can be framed as attempts to modify $\Delta E$ dynamics to move the system back toward higher-coherence basins.
\end{itemize}

\section{Experimental Paradigms}

To test and refine the ONTO$\Sigma$ interpretation of neural coherence and drift, we propose several experimental paradigms that explicitly manipulate impulses \(I_t\), track structural proxies for \(E_t\), and estimate coherence \(C_t\) across behavioural, neural, and subjective levels.

\subsection*{Wim Hof and controlled breathing studies}

Controlled breathing and cold exposure protocols, such as those used in the Wim Hof Method (WHM), offer a natural testbed:

\begin{itemize}[leftmargin=2em]
    \item \textbf{Impulses \(I_t\).} Cold exposure, breathing patterns, and interoceptive signals act as controlled impulses.
    \item \textbf{Structure proxies \(E_t\).} Longitudinal changes in:
    \begin{itemize}
        \item resting-state connectivity,
        \item task-evoked activation patterns,
        \item autonomic parameters (HRV, baroreflex sensitivity).
    \end{itemize}
    \item \textbf{Coherence measures \(C_t\).} Composite indices combining:
    \begin{itemize}
        \item performance on stress-challenge tasks,
        \item neural synchrony patterns during cold exposure,
        \item subjective ratings of control, meaning, and well-being.
    \end{itemize}
\end{itemize}

A UTAM-style hypothesis is that regular WHM practice induces $\Delta E$-like changes in brain and body structure that increase coherence \(C_t\) under cold and stress impulses, and that these changes can be quantified as movement into different attractor regimes in the volitional phase space.

\subsection*{Task-switching under hidden drift}

Another paradigm involves cognitive tasks whose rules or statistics drift over time without explicit cues:

\begin{itemize}[leftmargin=2em]
    \item Participants perform a task (e.g.\ categorization, decision-making) where the mapping from stimuli to correct responses changes slowly and covertly.
    \item \textbf{Impulses \(I_t\).} Trial-by-trial stimuli and feedback, with hidden drifts in underlying contingencies.
    \item \textbf{Structure proxies \(E_t\).} Trial-to-trial changes in inferred internal models, estimated via behavioural fits or neural decoding (e.g.\ changes in prefrontal representation).
    \item \textbf{Coherence \(C_t\).} Accuracy, reaction times, and measures of model fit, combined with neural indices of stable yet flexible network configuration.
\end{itemize}

In UTAM terms, such tasks explicitly probe the system's ability to maintain coherence under drift. Conditions in which adaptation is blocked or slowed (e.g.\ through cognitive load, fatigue, or pharmacological interventions) should show:

\begin{itemize}[leftmargin=2em]
    \item decreasing \(C_t\) and increasing \(D_t\) over time,
    \item characteristic shifts in connectivity (e.g.\ reduced fronto-parietal engagement),
    \item and subjective reports of confusion or effort.
\end{itemize}

\subsection*{Triangulating behavioural, neural, and subjective correlates}

A common structure across these paradigms is the triangulation of three levels:

\begin{description}[leftmargin=2em]
    \item[Behavioural.] Accuracy, response times, error patterns, performance under challenge.
    \item[Neural/physiological.] Functional connectivity, oscillatory coherence, autonomic markers, hormonal profiles.
    \item[Subjective.] Reports of clarity, control, stress, meaning, and self-coherence.
\end{description}

ONTO$\Sigma$ predicts that:

\begin{itemize}[leftmargin=2em]
    \item these three levels will covary in ways that track a latent coherence variable \(C_t\),
    \item and that interventions which modify $\Delta E$ (e.g.\ training, therapy, neurofeedback, breathing protocols) will move systems along identifiable trajectories in this joint space.
\end{itemize}

In this way, neuroscience and physiology provide a concrete arena in which the abstract notions of will, structure, action, coherence, and drift can be tested, refined, and potentially falsified. Subsequent chapters will extend this approach to AI, psychology, and social systems.


\chapter{AI, Machine Learning, and Robotics}

Artificial agents and robotic systems provide a particularly clear arena in which to test and develop the ONTO$\Sigma$ framework. Unlike brains and societies, their structures \(E_t\) and update rules \(G\) can be specified explicitly, modified at will, and evaluated under controlled conditions. At the same time, modern AI and robotics face exactly the challenge that UTAM is designed to diagnose: maintaining coherent behaviour in \emph{non-stationary} environments where naive, fixed controllers fail.

In this chapter we outline how I$^2$C and $\Delta E$ can be instantiated in AI and robotics, how coherence and drift can be measured in such systems, and how these ideas connect to contemporary concerns about robust, safe, and aligned AI.

\section{Non-Stationary Environments and Concept Drift}

Most classical AI and control designs implicitly assume that the environment is either stationary or slowly varying within a known envelope. Policies are trained on a fixed data distribution; controllers are tuned for a particular plant model. Yet real-world deployments almost always violate these assumptions:

\begin{itemize}[leftmargin=2em]
    \item sensor characteristics degrade or shift,
    \item user behaviour changes,
    \item hardware wears,
    \item the broader task context evolves.
\end{itemize}

In machine learning, this is known as \emph{concept drift}: the mapping from input features to labels or rewards changes over time. A model \(E_t\) that once performed well becomes misaligned with the actual impulse distribution \(I_t\); accuracy drops, and the system's behaviour becomes brittle or unsafe.

In control and robotics, the parallel phenomenon occurs when controllers are designed for a nominal model of the plant and environment, but actual dynamics drift outside the design envelope. Fixed-gain controllers can become unstable; policies learned in simulation fail on real hardware.

From the ONTO$\Sigma$ perspective, these are all instances of the same pattern:

\begin{itemize}[leftmargin=2em]
    \item The environment is non-stationary (\(dI/dt \neq 0\)).
    \item The structure-updater \(G\) is trivial or absent (\(\partial G / \partial I \approx 0\)).
    \item Coherence \(C_t\) between \(E_t\), \(I_t\), and \(A_t\) declines over time; drift \(D_t = 1 - C_t\) increases.
\end{itemize}

Static policies and fixed controllers thus realize a degenerate case of horizontal UTAM: they have an \(F(E_t, I_t)\) that generates actions but no significant $\Delta E$ process to keep \(E_t\) aligned with the evolving world. ONTO$\Sigma$ proposes that robust AI and robotics require explicit design at the structural level: controllers and agents must be endowed with $\Delta E$ dynamics that aim to preserve coherence, not just short-term performance.

\section{\texorpdfstring{$\Delta E$}{ΔE} Agents vs.\ Classical Controllers}

We now sketch how $\Delta E$ can be implemented in artificial agents and robotic controllers, and how coherence measures can be integrated into existing learning frameworks.

\subsection*{Designing $\Delta E$-based controllers in robotics}

Consider a robot interacting with an environment whose dynamics drift over time (changing terrain, payload, friction, or contact conditions). We can interpret:

\begin{itemize}[leftmargin=2em]
    \item the robot's control policy or internal model as its structure \(E_t\),
    \item the sensor readings as impulses \(I_t\),
    \item the motor commands as actions \(A_t\).
\end{itemize}

A $\Delta E$-based controller would then:

\begin{enumerate}[leftmargin=2em]
    \item implement an I$^2$C cycle at each control step:
    \[
        Z_t = \mathrm{Enc}(I_t, E_t), \qquad
        A_t = \arg\max_{A} C(Z_t, A, E_t),
    \]
    where \(C\) encodes criteria such as stability, safety, and task progress;
    \item update its internal structure \(E_t\) according to a coherence-preserving rule:
    \[
        E_{t+1} = E_t + \eta \, \nabla_E C(E_t, I_t, A_t),
    \]
    or a more domain-specific $\Delta E$ variant (e.g.\ adaptive parameter estimation, online system identification, meta-learning).
\end{enumerate}

Unlike classical adaptive control, which often focuses on maintaining tracking performance for a given reference signal, $\Delta E$-based controllers can be designed to preserve a richer notion of coherence that includes:

\begin{itemize}[leftmargin=2em]
    \item structural integrity (avoiding configurations that damage hardware),
    \item safety margins (staying away from dangerous state regions),
    \item and long-term viability (limiting wear, energy consumption).
\end{itemize}

In UTAM terms, such a controller is an artificial manifestation of will that explicitly protects its own structural capacity to act.

\subsection*{Integrating coherence into reinforcement learning}

In reinforcement learning, an agent seeks to maximize expected cumulative reward. From a UTAM standpoint, this makes reward the primary driver of behaviour, while coherence is at best an indirect concern. To move toward $\Delta E$ agents, we can:

\begin{itemize}[leftmargin=2em]
    \item introduce an explicit coherence term \(C_t\) into the objective, so that policies are trained to maximise a combination of external reward and internal coherence;
    \item treat parts of the policy or value function parameters as \(E_t\) and update them using $\Delta E$ rules that are sensitive not only to reward gradients but also to changes in the impulse distribution;
    \item maintain a separate coherence-monitoring module that evaluates how well the current structure matches ongoing data, and triggers structural updates (e.g.\ representation learning, architecture adaptation) when drift is detected.
\end{itemize}

Concretely, this might mean:

\begin{align*}
    A_t &= \pi_{\theta_t}(I_t) \quad \text{(I$^2$C action)}, \\
    \theta_{t+1} &= \theta_t + \eta_1 \nabla_{\theta} \mathbb{E}[R_t] 
                   + \eta_2 \nabla_{\theta} C(\theta_t, I_t, A_t),
\end{align*}
where \(\theta_t\) are the agent's parameters, \(\mathbb{E}[R_t]\) is expected reward, and \(C\) is a coherence functional (e.g.\ prediction accuracy, uncertainty calibration, safety constraint satisfaction).

In this way, RL agents become \(\Delta E\) agents: they do not only learn to get reward; they learn to maintain and improve the coherence of their own structures in the face of drift.

\section{Benchmark Proposals}

To empirically evaluate the benefits of I$^2$C + $\Delta E$ designs over static or purely reward-driven approaches, we propose a family of benchmarks that explicitly incorporate environmental drift and structural adaptation.

\subsection*{Drifting terrain tasks in robotics}

A natural class of benchmarks involves robots operating on terrains whose properties change over time:

\begin{itemize}[leftmargin=2em]
    \item legged robots moving from flat ground to slopes, loose soil, or ice;
    \item manipulators handling objects whose mass distribution or friction properties drift;
    \item aerial robots encountering changing wind patterns or sensor degradation.
\end{itemize}

In each case, one can compare:

\begin{itemize}[leftmargin=2em]
    \item \textbf{Fixed controllers} tuned for nominal conditions;
    \item \textbf{Standard adaptive controllers} (e.g.\ gain scheduling, model reference adaptive control);
    \item \textbf{$\Delta E$ controllers} that explicitly monitor coherence and update their internal structure accordingly.
\end{itemize}

Performance metrics would include:

\begin{itemize}[leftmargin=2em]
    \item task performance (distance travelled, error, energy use),
    \item coherence measures (prediction error, stability margins, safety incidents),
    \item adaptation time (how quickly coherence is restored after a drift event).
\end{itemize}

\subsection*{Non-stationary RL benchmarks}

In the RL setting, we can design benchmarks where the environment's transition and reward functions drift over time:

\begin{itemize}[leftmargin=2em]
    \item \textbf{Piecewise-stationary MDPs}, where the environment switches between distinct regimes (e.g.\ different dynamics or reward structures) at unknown times.
    \item \textbf{Gradual drift environments}, where parameters of the dynamics or reward functions change slowly and continuously.
\end{itemize}

Agents can then be evaluated on:

\begin{itemize}[leftmargin=2em]
    \item cumulative reward under drift,
    \item coherence metrics (e.g.\ calibration of predictive models, constraint satisfaction),
    \item robustness (graceful degradation when overwhelmed),
    \item and adaptation latency (time from a drift event to recovery of adequate coherence).
\end{itemize}

ONTO$\Sigma$ predicts that agents with explicit $\Delta E$ mechanisms---those that treat their internal structure as an object of continual coherence-preserving adaptation---will exhibit superior robustness and more interpretable failure modes than agents that rely solely on static policies or standard online RL updates tuned purely to reward.

\section{Toward Volitional Agents}

Finally, we consider how ONTO$\Sigma$ reframes the long-term goal of AI not just as building intelligent optimisers, but as designing systems that embody a form of \emph{artificial volition}: agents that maintain their own structural integrity as manifestations of directedness.

\subsection*{Beyond reward maximization}

Traditional AI benchmarks and designs prioritise performance on externally defined tasks: games won, errors minimized, goals reached. From a UTAM vantage point, this corresponds to focusing on \emph{particular} actions and outcomes while neglecting the \emph{ongoing viability} of the structures that make action possible.

A \emph{volitional agent} in the ONTO$\Sigma$ sense would:

\begin{itemize}[leftmargin=2em]
    \item treat the preservation and enhancement of coherence \(C_t\) as a primary objective,
    \item update its internal structure \(E_t\) via $\Delta E$ rules that reflect both external feedback and internal consistency constraints,
    \item avoid pursuing short-term reward in ways that erode its own capacity for coherent action.
\end{itemize}

Such an agent is not merely an optimiser of a fixed utility function; it is a system that \emph{cares} (in a formal, structural sense) about remaining a viable, coherent expression of directedness across changing contexts.

\subsection*{Links to safe and aligned AI}

Concerns about AI safety and alignment centre on the possibility that powerful agents might:

\begin{itemize}[leftmargin=2em]
    \item pursue goals in ways that damage their environment or themselves,
    \item generalize poorly under distribution shift,
    \item or become brittle, unpredictable, or uncontrollable outside training regimes.
\end{itemize}

ONTO$\Sigma$ offers a complementary perspective on these issues:

\begin{itemize}[leftmargin=2em]
    \item Agents that maintain high coherence \(C_t\) with respect to \emph{both} their environment and their own structural integrity are less likely to exhibit catastrophic misbehaviour under drift.
    \item $\Delta E$ dynamics can be designed to incorporate safety constraints and alignment criteria directly into the coherence functional \(C\), so that unsafe or misaligned structural changes are penalized at the level of adaptation, not just at the level of action outcomes.
    \item Monitoring coherence and drift provides a natural set of \emph{early-warning signals}: increasing \(D_t\) can flag that an agent is entering regimes where its learned policies no longer apply, prompting human oversight or structural intervention.
\end{itemize}

In this way, the notion of \emph{volitional agents} does not require anthropomorphizing AI systems. It simply means designing agents whose internal dynamics approximate the UTAM pattern:

\[
    \Will \to E_t \to A_t, \quad
    A_t = F(E_t, I_t), \quad
    E_{t+1} = G(E_t, A_t, I_t, \Theta),
\]
with \(G\) chosen so that coherence is preserved and drift is managed rather than ignored. Whether or not we treat such agents as having ``real'' will in a philosophical sense, ONTO$\Sigma$ argues that this volitional architecture---structure-preserving, coherence-oriented adaptation under drift—is a necessary ingredient of robust, safe, and truly \emph{aligned} artificial systems.


\chapter{Psychology, Self, and Meaning}

ONTO$\Sigma$ is not merely a theory about abstract systems; it is also a proposal about how to understand human psychology, selfhood, and meaning. In this chapter we reinterpret familiar psychological notions in UTAM terms, treating the \emph{self} as a coherence structure, \emph{meaning} as a long-term coherence regime, and \emph{burnout} or \emph{depersonalization} as drift phenomena. We also outline how narrative and therapeutic practices can be seen as forms of $\Delta E$ training.

\section{Self as Coherence Structure}

In everyday language, the \emph{self} is often treated as a thing: an inner entity that has experiences and makes decisions. ONTO$\Sigma$ instead treats the self as a \emph{pattern of coherence} within the ongoing dynamics of will, structure, and action.

Let us introduce a time-indexed self-structure \(\text{Self}_t\), understood as the current organisation of:

\begin{itemize}[leftmargin=2em]
    \item attentional patterns (what is being noticed and foregrounded),
    \item interpretive schemas (how impulses \(I_t\) are made sense of),
    \item action policies (how responses \(A_t\) are selected),
    \item and meta-constraints (values, commitments, identity claims).
\end{itemize}

Formally, we can think of \(\text{Self}_t\) as a functional of the system's structure and current focus:
\[
    \text{Self}_t \approx \mathcal{S}(E_t, \alpha_t),
\]
where \(\alpha_t\) encodes attentional allocation (which components of \(I_t\) and \(E_t\) are emphasized). The self at time \(t\) is then the \emph{particular way} in which \Will{} is being organized into a coherent pattern of attending, interpreting, and acting.

In UTAM terms:

\begin{itemize}[leftmargin=2em]
    \item \(\Will\) provides the basic directedness.
    \item \(E_t\) provides the structural repertoire of interpretations and actions.
    \item \(\text{Self}_t\) is the \emph{current configuration} of that repertoire---a local, coherent alignment of will, attention, interpretation, and action.
\end{itemize}

A stable sense of self then corresponds to a relatively stable attractor in the structural manifold \(\mathcal{M}\), where the system returns to a characteristic pattern of coherence after perturbations.

\section{Meaning as Long-Term Coherence Regime}

If the self is a local coherence structure, \emph{meaning} can be understood as a \emph{long-term coherence regime} under drift. Intuitively, a life feels meaningful when:

\begin{itemize}[leftmargin=2em]
    \item actions over extended periods form a coherent pattern,
    \item this pattern remains robust in the face of changing circumstances,
    \item and it integrates local episodes into a recognisable trajectory.
\end{itemize}

In the language of volitional phase space, meaning corresponds to a trajectory
\[
    \gamma(t) = (E_t, I_t)
\]
that:

\begin{itemize}[leftmargin=2em]
    \item stays within or near high-coherence regions of \(C(E, I)\) over long timescales,
    \item and does so despite non-trivial drift in \(I_t\) (changing roles, crises, development).
\end{itemize}

We can thus say:

\begin{quote}
A life is \emph{meaningful}, in the ONTO$\Sigma$ sense, to the extent that it realizes a stable, high-coherence regime of \(\Will\)–\(\Struct_t\)–\(\Act_t\) across the drifts of circumstance.
\end{quote}

Subjectively, this appears as a sense that one's actions ``hang together'', that there is a through-line connecting past, present, and future, and that demands from the environment can be integrated rather than merely endured. Objectively, it may be reflected in:

\begin{itemize}[leftmargin=2em]
    \item consistent value-expressive behaviour,
    \item resilient yet flexible coping strategies,
    \item and long-term projects that maintain or deepen coherence rather than eroding it.
\end{itemize}

\section{Burnout and Depersonalization as Drift}

Within this framework, \emph{burnout} and \emph{depersonalization} are not just emotional labels; they are names for particular \emph{drift regimes}.

\subsection*{Burnout as chronic misalignment}

Burnout can be characterized by:

\begin{itemize}[leftmargin=2em]
    \item persistent low \(C_t\) despite ongoing effort,
    \item a sense that actions no longer restore coherence but merely postpone breakdown,
    \item and a weakening of $\Delta E$ dynamics (the feeling that ``nothing I change really helps'').
\end{itemize}

In volitional phase space, the system has been pulled into a basin where:

\begin{itemize}[leftmargin=2em]
    \item the current structure \(E_t\) is mismatched to the dominant impulses \(I_t\),
    \item the gradient \(\nabla_E C\) is shallow, offering little guidance for adaptation,
    \item and repeated attempts to adjust within this basin lead to negligible increases in coherence.
\end{itemize}

The phenomenology of exhaustion, cynicism, and reduced efficacy can thus be seen as first-person access to a state of high drift \(D_t\) and low, flat coherence landscape around the current self-structure.

\subsection*{Depersonalization as loss of self-coherence}

\emph{Depersonalization}---the feeling of being unreal, detached from one's own actions or experiences---can be interpreted as a breakdown in the alignment that normally defines \(\text{Self}_t\):

\begin{itemize}[leftmargin=2em]
    \item attentional focus no longer tracks what matters to the person,
    \item interpretations feel alien or arbitrary,
    \item actions are experienced as if they ``belonged to someone else''.
\end{itemize}

Formally, this corresponds to a regime in which the mapping
\[
    \mathcal{S}(E_t, \alpha_t) \mapsto \text{Self}_t
\]
becomes unstable or fragmented: different parts of \(E_t\) and \(\alpha_t\) no longer cohere into a single, integrated pattern. This is a particularly vivid example of drift: \(\Will\) is still being expressed through structures and actions, but the system's own capacity to ``own'' this expression as a unified self has degraded.

\section{Narrative Coherence and Well-Being}

One of the most accessible windows into long-term coherence regimes is narrative. Life stories---the accounts individuals give of their past, present, and anticipated future---can be viewed as \emph{symbolic projections} of trajectories in volitional phase space.

\subsection*{Life stories as trajectories}

Given a trajectory \(\gamma(t) = (E_t, I_t)\), a narrative is a higher-level structure \(N\) that:

\begin{itemize}[leftmargin=2em]
    \item selects certain events and actions as salient,
    \item connects them via causal and thematic links,
    \item and assigns them roles in a larger pattern (e.g.\ growth, redemption, tragedy).
\end{itemize}

A \emph{narratively coherent} life story is one in which:

\begin{itemize}[leftmargin=2em]
    \item the selected events form an intelligible progression,
    \item apparent contradictions can be reconciled or integrated,
    \item and the implied future maintains or restores coherence rather than shattering it.
\end{itemize}

In ONTO$\Sigma$ terms, such a narrative reflects an underlying trajectory that remains, on balance, in high-coherence regions of \(C(E, I)\), or at least moves toward them after perturbations.

\subsection*{Empirical study designs and metrics}

This perspective suggests concrete empirical programmes:

\begin{itemize}[leftmargin=2em]
    \item \textbf{Narrative analysis.} Collect life stories and code them for:
    \begin{itemize}
        \item temporal coherence (clear ordering, causal links),
        \item thematic coherence (consistent identity themes, values),
        \item integrative capacity (ability to incorporate adversity).
    \end{itemize}
    \item \textbf{Well-being measures.} Administer standard scales for:
    \begin{itemize}
        \item life satisfaction and meaning in life,
        \item symptoms of burnout, depersonalization, depression, and anxiety,
        \item resilience and post-traumatic growth.
    \end{itemize}
    \item \textbf{Longitudinal tracking.} Follow individuals over time to see how changes in narrative coherence relate to:
    \begin{itemize}
        \item changes in behaviour (career shifts, relationship patterns),
        \item changes in subjective coherence (self-report \(C_t\)-like indices),
        \item and changes in physiological markers (e.g.\ stress-related measures).
    \end{itemize}
\end{itemize}

The UTAM prediction is that higher narrative coherence will correlate with higher long-term coherence \(C_t\) in behaviour and physiology, and that shifts in narrative structure can function as a form of $\Delta E$ update---a reorganization of \(E_t\) that changes how future impulses \(I_t\) are interpreted and acted upon.

\section{Therapeutic Interventions as \texorpdfstring{$\Delta E$}{ΔE} Training}

Finally, ONTO$\Sigma$ offers a unifying way to think about psychotherapy and related interventions: as \emph{training in structural update}, rather than merely as goal-setting or symptom management.

\subsection*{Reframing therapy as structural update}

Most therapeutic modalities involve, in one way or another:

\begin{itemize}[leftmargin=2em]
    \item identifying patterns in how a person interprets impulses and selects actions,
    \item highlighting where these patterns produce incoherence or unnecessary drift,
    \item and practising alternative interpretations and responses that lead to greater coherence.
\end{itemize}

In UTAM language, this is precisely $\Delta E$ work:

\begin{itemize}[leftmargin=2em]
    \item \(\text{Impulses } I_t\): life events, emotions, thoughts, bodily sensations.
    \item \(\text{Current structure } E_t\): beliefs, schemas, habits, attachment patterns.
    \item \(\text{Actions } A_t\): behavioural responses, cognitive strategies, interpersonal moves.
    \item \(\Delta E\): explicit, guided updates to \(E_t\) via new interpretations, experiments, and practices.
\end{itemize}

Therapy is not just about choosing different goals; it is about changing what counts as a coherent response in the first place.

\subsection*{Practical examples and possible protocols}

Different therapeutic approaches can be reinterpreted as emphasising different aspects of $\Delta E$:

\begin{itemize}[leftmargin=2em]
    \item \textbf{Cognitive–behavioural therapy (CBT).} Identifies distorted interpretations and maladaptive action patterns; experiments with alternative thoughts and behaviours to increase coherence between beliefs, evidence, and outcomes. In UTAM terms: targeted $\Delta E$ updates to cognitive substructures of \(E_t\).
    \item \textbf{Acceptance and commitment therapy (ACT).} Clarifies values (long-term coherence criteria), fosters acceptance of unchangeable impulses \(I_t\), and trains committed action aligned with values. Here, \(\Delta E\) involves reshaping the coherence functional \(C\) itself by making values explicit and behaviourally grounded.
    \item \textbf{Trauma-focused therapies.} Help re-encode traumatic impulses \(I_t\) and their traces in \(E_t\) so that they no longer dominate coherence evaluations. Techniques such as exposure, EMDR, or somatic work can be seen as controlled manipulations of \(I_t\) and \(A_t\) designed to move the system from a trauma-shaped basin to a healthier coherence regime.
\end{itemize}

Future protocols inspired by ONTO$\Sigma$ could explicitly:

\begin{itemize}[leftmargin=2em]
    \item measure subjective and behavioural proxies of \(C_t\) and \(D_t\) over the course of therapy,
    \item frame interventions as experiments in the volitional phase space (trying out new \(E_t\) in different \(I_t\) contexts),
    \item and train clients in meta-structural skills: the ability to recognize when drift is occurring and to initiate their own $\Delta E$ processes (self-guided structural updates).
\end{itemize}

In this way, psychology and psychotherapy become domains in which UTAM's abstract notions of will, structure, action, coherence, drift, and $\Delta E$ are given concrete, human shape. The self is no longer a mysterious substance but a dynamic coherence structure; meaning is not an ineffable property but a long-term pattern of high-coherence trajectories; and healing is the art of guiding systems back into regimes where will can once again be coherently embodied and enacted.

\chapter{Social, Institutional, and Civilizational Dynamics}

Up to this point, we have applied UTAM to individual systems: brains, agents, robots. But will (\(\Will\)) also manifests at larger scales, through collective structures that shape the behaviour of many individuals at once. Laws, norms, infrastructures, markets, and cultural narratives are all ways in which directedness becomes organized beyond the single organism.

In this chapter we extend the ONTO$\Sigma$ vocabulary to social, institutional, and civilizational dynamics. We treat collective structures as instances of \(E_t\), examine how institutional coherence and drift arise, and outline how agent-based models with I$^2$C + \(\Delta E\) can be used to explore attractors and phase transitions in social systems.

\section{Collective Structures as \texorpdfstring{$E_t$}{E\_t}}

At the social level, structure \(E_t\) is no longer just the internal model of a single agent. It includes:

\begin{itemize}[leftmargin=2em]
    \item \textbf{Laws and formal rules:} constitutions, regulations, contracts, procedures.
    \item \textbf{Norms and informal expectations:} customs, moral codes, professional standards.
    \item \textbf{Infrastructures:} physical networks (roads, power grids), information systems (media, platforms), economic arrangements (currencies, markets).
\end{itemize}

All of these can be understood as \emph{structural manifestations of will at scale}: coordinated patterns through which collective directedness is embodied and channelled. In UTAM notation, we can speak of:

\[
    E_t^{\mathrm{macro}} \in \mathcal{E}(\Will)
\]

as a macro-structure that constrains and enables the actions of many micro-level agents. Individual agents still have their own structures \(E_t^{(i)}\), but these are embedded in and co-shaped by the institutional \(E_t^{\mathrm{macro}}\).

From an ONTO$\Sigma$ perspective:

\begin{itemize}[leftmargin=2em]
    \item \(\Will\) is expressed not only in individual intentions but also in the way collective life is organized.
    \item Institutions, norms, and infrastructures are ways in which directedness is \emph{stabilized} and \emph{amplified} across many agents and over long times.
\end{itemize}

The question of social health, then, is a question of how coherent these collective structures are with respect to the underlying conditions, impulses, and actions they are meant to organize.

\section{Institutional Coherence and Drift}

Just as individuals can be more or less coherent, institutions can exhibit varying degrees of \emph{institutional coherence}. Informally, an institution is coherent when:

\begin{itemize}[leftmargin=2em]
    \item its formal rules and informal norms \(E_t^{\mathrm{macro}}\) fit the actual environment (economic, technological, ecological, demographic),
    \item the actions it generates through its procedures and incentives \(A_t^{\mathrm{macro}}\) effectively realize its stated purposes,
    \item and its internal feedback mechanisms allow it to adapt when conditions change.
\end{itemize}

We can thus define a macro-level coherence \(C_t^{\mathrm{inst}}\) that measures how well an institution's structure fits:

\begin{itemize}[leftmargin=2em]
    \item the impulses it faces (public needs, resource constraints, shocks),
    \item the actions it takes (policies, enforcement, investment),
    \item and the downstream effects on the populations and systems it governs.
\end{itemize}

\subsection*{Adaptive vs.\ rigid institutions}

Adaptive institutions approximate \(\Delta E\) processes at the structural level:

\begin{itemize}[leftmargin=2em]
    \item They monitor their own performance (coherence proxies: trust, effectiveness, legitimacy).
    \item They incorporate feedback from impulses \(I_t^{\mathrm{macro}}\) (public opinion, economic indicators, scientific knowledge).
    \item They update their structures \(E_t^{\mathrm{macro}}\)---laws, norms, procedures---to improve coherence.
\end{itemize}

Examples of $\Delta E$-like mechanisms include:

\begin{itemize}[leftmargin=2em]
    \item democratic elections and institutional reforms,
    \item independent oversight bodies and audits,
    \item adaptive regulation (sunset clauses, iterative policy design),
    \item professional self-regulation and evolving standards.
\end{itemize}

Rigid institutions, by contrast, behave as if \(\partial G / \partial I \approx 0\): their structural update mechanisms are frozen or purely symbolic. When the environment drifts, their internal dynamics continue as before, leading to increasing misalignment:

\begin{itemize}[leftmargin=2em]
    \item policies that systematically fail to address actual problems,
    \item norms that are widely violated but rarely revised,
    \item infrastructures that no longer match usage patterns or technological baselines.
\end{itemize}

In UTAM terms, such institutions are in a state of \emph{institutional drift}: \(D_t^{\mathrm{inst}} = 1 - C_t^{\mathrm{inst}}\) grows over time, and small shocks can trigger outsized failures because the structural fit has been eroded.

\subsection*{Feedback mechanisms as \texorpdfstring{$\Delta E$}{ΔE}-like processes}

We can explicitly interpret various societal feedback mechanisms as macro-level $\Delta E$ processes:

\begin{itemize}[leftmargin=2em]
    \item \textbf{Elections and public deliberation} update the composition and priorities of political structures when coherence between citizens' needs and institutional actions declines.
    \item \textbf{Market signals} (prices, investment flows) update economic structures, reallocating resources away from incoherent arrangements.
    \item \textbf{Social movements and cultural criticism} highlight incoherence between professed values and actual practices, pushing norms and laws toward new coherence regimes.
\end{itemize}

When these feedback loops are healthy, they act as distributed $\Delta E$: they adjust the institutional \(E_t^{\mathrm{macro}}\) in response to drift in \(I_t^{\mathrm{macro}}\). When they are captured, suppressed, or drowned in noise, institutional \(E_t^{\mathrm{macro}}\) becomes rigid, and drift accelerates.

\section{Civilizational Drift and Collapse}

At the broadest scale, we can talk about \emph{civilizational} structures: large-scale patterns of production, governance, culture, and technology that define an era. In ONTO$\Sigma$ terms, a civilization corresponds to a high-level \(E_t^{\mathrm{civ}}\) that:

\begin{itemize}[leftmargin=2em]
    \item encodes dominant institutions, infrastructures, and narratives,
    \item organizes the actions of millions or billions of individuals,
    \item and mediates the relationship between \Will{} and planetary-scale impulses \(I_t^{\mathrm{civ}}\) (ecological conditions, resource availability, technological capabilities).
\end{itemize}

\subsection*{Stagnation and overreaction}

Civilizational drift can manifest both as \emph{stagnation} and as \emph{overreaction}:

\begin{itemize}[leftmargin=2em]
    \item In stagnation, \(E_t^{\mathrm{civ}}\) changes very slowly while \(I_t^{\mathrm{civ}}\) drifts significantly (e.g.\ ecological degradation, demographic shifts). Coherence \(C_t^{\mathrm{civ}}\) erodes quietly; institutions continue to operate according to scripts that no longer fit realities.
    \item In overreaction, structural updates are frequent but poorly grounded: reforms, revolutions, or technological deployments that are not guided by coherent $\Delta E$ principles. The result is volatility rather than restored coherence; the system oscillates between mismatched structures rather than converging on a stable regime.
\end{itemize}

Both patterns reflect failures in macro-scale I$^2$C and $\Delta E$: either impulses are not properly interpreted and integrated, or structural updates are decoupled from genuine coherence gradients.

\subsection*{Loss of structural fit and collapse}

In extreme cases, civilizational drift can lead to partial or total collapse:

\begin{itemize}[leftmargin=2em]
    \item key infrastructures fail (food, water, energy, information),
    \item trust in institutions disintegrates,
    \item and large-scale coordination breaks down.
\end{itemize}

From the UTAM viewpoint, collapse is the endpoint of a trajectory in which:

\begin{itemize}[leftmargin=2em]
    \item \(D_t^{\mathrm{civ}}\) has been increasing for a long time,
    \item \(\Delta E\) mechanisms were either absent, captured, or misapplied,
    \item and the coherence landscape \(C(E^{\mathrm{civ}}, I^{\mathrm{civ}})\) has been deformed by underlying conditions (e.g.\ resource limits) in ways that existing structures cannot track.
\end{itemize}

Importantly, ONTO$\Sigma$ does not treat collapse as fate. It is the consequence of specific failures in macro-level I$^2$C and $\Delta E$ that can, in principle, be mitigated: by strengthening feedback loops, expanding the accessible region of \(\mathcal{E}(\Will)\) through innovations, and designing institutions that explicitly monitor and preserve coherence under long-term drift.

\section{Agent-Based Models with UTAM Dynamics}

To move from qualitative narratives to quantitative exploration, we can build \emph{agent-based models} (ABMs) in which UTAM dynamics are implemented at multiple scales.

\subsection*{Populations of UTAM agents}

In a UTAM-style ABM:

\begin{itemize}[leftmargin=2em]
    \item Each agent \(i\) has its own structure \(E_t^{(i)}\) and impulses \(I_t^{(i)}\), and follows I$^2$C + $\Delta E$ dynamics:
    \[
        A_t^{(i)} = F(E_t^{(i)}, I_t^{(i)}), \qquad
        E_{t+1}^{(i)} = G(E_t^{(i)}, A_t^{(i)}, I_t^{(i)}, \Theta^{(i)}).
    \]
    \item Agents interact with each other and with shared environments, so that:
    \begin{itemize}
        \item each agent's \(I_t^{(i)}\) depends on others' actions \(A_t^{(j)}\),
        \item aggregate patterns of \(E_t^{(i)}\) and \(A_t^{(i)}\) give rise to emergent macro-structures \(E_t^{\mathrm{macro}}\).
    \end{itemize}
\end{itemize}

We can then study how different designs of \(F\), \(G\), and environmental dynamics affect:

\begin{itemize}[leftmargin=2em]
    \item the emergence of institutions (norms, hierarchies, markets),
    \item the resilience or fragility of these institutions under shocks,
    \item the distribution of coherence and drift across agents.
\end{itemize}

\subsection*{Exploring attractors and phase transitions in social systems}

Within such models, the volitional phase space extends to include both individual and collective structures:

\[
    \mathcal{V}_{\mathrm{soc}} \approx 
        \bigl(\mathcal{M}_{\mathrm{ind}}^{N} \times \mathcal{M}_{\mathrm{macro}}\bigr)
        \times \mathcal{I}_{\mathrm{soc}},
\]
where:

\begin{itemize}[leftmargin=2em]
    \item \(\mathcal{M}_{\mathrm{ind}}\) is the structural manifold for individual agents,
    \item \(\mathcal{M}_{\mathrm{macro}}\) for institutions,
    \item \(\mathcal{I}_{\mathrm{soc}}\) the space of social impulses (shocks, resources, information flows).
\end{itemize}

Simulations can then be used to explore:

\begin{itemize}[leftmargin=2em]
    \item \textbf{Attractors:} stable patterns such as cooperative equilibria, polarization regimes, or hierarchical orders that act as coherence basins in \(\mathcal{V}_{\mathrm{soc}}\).
    \item \textbf{Phase transitions:} abrupt shifts from one regime to another as parameters (e.g.\ resource availability, communication bandwidth, adaptation rates \(\eta\)) cross critical thresholds.
    \item \textbf{Meta-structural interventions:} policies or innovations that change the form of \(G\) (how institutions learn) or the definition of \(C\) (what counts as coherent), and their impact on long-term trajectories.
\end{itemize}

In this way, ONTO$\Sigma$ provides not only a philosophical lens on social and civilizational dynamics, but also a formal template for constructing models that can generate testable hypotheses about how collective coherence is maintained, lost, or transformed over time.


% ================================================
% Part IV – Formal and Philosophical Deepening
% ================================================

\part{Formal and Philosophical Deepening}

\chapter{Formal Foundations}

In the preceding parts we have introduced UTAM at three levels:

\begin{itemize}[leftmargin=2em]
    \item an \emph{ontological level}, where \(\Will \to \Struct \to \Act\) expresses the priority of will, structure, and action;
    \item a \emph{dynamical level}, where I$^2$C and \(\Delta E\) define horizontal flows on the structural manifold under impulses;
    \item a \emph{geometric level}, where the volitional phase space and coherence landscape provide intuition for attractors, drift, and meta-structural dynamics.
\end{itemize}

In this chapter we sketch a more formal foundation for these ideas. The aim is not yet to fix a single definitive mathematics for UTAM, but to outline how coherence, $\Delta E$ dynamics, and the vertical/horizontal distinction can be rendered in rigorous terms using existing tools from information theory, control theory, and category theory.

\section{Rigorous Definitions of Coherence}

Coherence \(C_t\) has so far been treated as an abstract scalar in \([0,1]\) that measures how well a given structure \(E_t\) handles impulses \(I_t\) by producing actions \(A_t\). To move toward formalization, we consider several families of candidate definitions.

\subsection*{Information-theoretic coherence}

Suppose we model the environment and system as random variables on a probability space. Let \(X_t\) denote impulses, \(Y_t\) actions, and \(Z_t\) internal states determined by \(E_t\). Information theory then offers multiple ways to define coherence:

\begin{itemize}[leftmargin=2em]
    \item \textbf{Predictive coherence.} Given a generative model \(p_{E_t}(X_{t+1} \mid \text{history})\), we can define coherence as a monotone transform of negative log-loss or KL divergence between predicted and actual impulse distributions:
    \[
        C_t^{\mathrm{pred}} = f\!\left(- D_{\mathrm{KL}}\bigl(
            p_{\mathrm{env}}(X_{t+1} \mid \text{history})
            \,\Vert\, p_{E_t}(X_{t+1} \mid \text{history})
        \bigr)\right).
    \]
    \item \textbf{Mutual-information coherence.} Coherence can track how much information about relevant aspects of the environment is retained and used in action:
    \[
        C_t^{\mathrm{MI}} = f\!\bigl(I(X_t; Y_t \mid E_t)\bigr),
    \]
    where \(I(\cdot;\cdot)\) is mutual information.
    \item \textbf{Free-energy coherence.} In variational formulations, coherence can be related to negative variational free energy:
    \[
        C_t^{\mathrm{FEP}} = f\!\bigl(-F(E_t; \text{data}_t)\bigr),
    \]
    where \(F\) measures a trade-off between accuracy and complexity for a generative model.
\end{itemize}

In each case, \(f\) is a monotone mapping that normalizes the raw quantity to \([0,1]\). The choice among these options depends on the domain and on which aspects of coherence are most salient (prediction, compression, controllability, etc.).

\subsection*{Control-theoretic coherence}

In control theory, coherence can be associated with stability, robustness, and performance relative to constraints:

\begin{itemize}[leftmargin=2em]
    \item \textbf{Stability margins.} For a closed-loop system with structure \(E_t\) (controller + plant model), we can let
    \[
        C_t^{\mathrm{stab}} = f(\text{gain margin}, \text{phase margin}, \text{Lyapunov exponents}, \dots),
    \]
    with higher coherence corresponding to stronger stability properties around relevant operating points.
    \item \textbf{Robust performance.} Coherence can quantify how well the system maintains output within viable bounds under disturbances and model uncertainty:
    \[
        C_t^{\mathrm{rob}} = f\bigl(\sup_{\Delta \in \mathcal{U}} \text{perf}(E_t,\Delta)\bigr),
    \]
    where \(\mathcal{U}\) is a set of admissible perturbations.
\end{itemize}

Here, coherence is explicitly about \emph{remaining viable} under drift, rather than hitting a single setpoint optimally.

\subsection*{Hybrid and phenomenological constraints}

In psychological and phenomenological contexts, coherence cannot be reduced to a single information-theoretic or control-theoretic metric. Instead, we consider hybrid measures:

\begin{itemize}[leftmargin=2em]
    \item composite indices that combine predictive accuracy, stability, and constraint satisfaction,
    \item plus phenomenological constraints such as:
    \begin{itemize}
        \item internal consistency of values and beliefs,
        \item felt sense of integrity or ``being oneself'',
        \item subjective meaning and agency.
    \end{itemize}
\end{itemize}

Formally, one can introduce a vector-valued coherence measure
\[
    \mathbf{C}_t = (C_t^{(1)}, C_t^{(2)}, \dots, C_t^{(k)}),
\]
and define scalar coherence as a monotone aggregation:
\[
    C_t = \Phi(\mathbf{C}_t),
\]
where \(\Phi\) encodes how different coherence aspects are weighted or traded off. This accommodates the fact that, for complex systems, coherence is multi-dimensional.

The key UTAM requirement is structural, not numerical: \(C_t\) must track how well \(\Will\)–\(E_t\)–\(A_t\) remain aligned under \(I_t\), and $\Delta E$ should tend to move the system toward higher \(C_t\) in expectation.

\section{Properties of \texorpdfstring{$\Delta E$}{ΔE} Dynamics}

We now turn to the dynamics of $\Delta E$ itself. Given a coherence functional \(C(E, I)\), we consider update rules of the general form:
\[
    E_{t+1} = E_t + \eta_t \, H_t(E_t, I_t, A_t),
\]
where \(H_t\) is some update direction (gradient-like or otherwise), and \(\eta_t > 0\) is a step size.

\subsection*{Convergence and local optima}

In the simplest case, \(H_t = \nabla_E C(E_t, I_t, A_t)\) and \(I_t\) is fixed. Then $\Delta E$ reduces to gradient ascent on a static landscape:
\[
    E_{t+1} = E_t + \eta \nabla_E C(E_t, I),
\]
or, in continuous time,
\[
    \frac{dE}{dt} = \eta \nabla_E C(E, I).
\]

Under standard regularity assumptions (e.g.\ smoothness of \(C\), small enough \(\eta\)), classical results apply:

\begin{itemize}[leftmargin=2em]
    \item trajectories converge to \emph{critical points} of \(C\) (where \(\nabla_E C = 0\));
    \item stable equilibria correspond to local maxima of \(C\) in \(\mathcal{M}\);
    \item the basin of attraction of a given local maximum defines a coherence attractor.
\end{itemize}

From the UTAM perspective, this means that $\Delta E$ dynamics will typically settle into locally coherent regimes, not necessarily globally optimal ones. This aligns with the phenomenology of habit formation and institutional path dependence: systems often find and reinforce locally coherent patterns that are difficult to escape even if higher coherence is available elsewhere in \(\mathcal{M}\).

\subsection*{Stability analysis}

Given a candidate equilibrium \(E^\ast\) with fixed impulse statistics \(I\), one can perform a standard stability analysis:

\begin{itemize}[leftmargin=2em]
    \item Linearize the dynamics around \(E^\ast\) to obtain a Jacobian \(J = \partial (\eta \nabla_E C)/\partial E \vert_{E^\ast}\).
    \item Analyse the eigenvalues of \(J\): negative real parts (in gradient \emph{descent}) correspond to stable minima; for gradient \emph{ascent} on \(C\), stable maxima require eigenvalues with negative real part for \(-J\).
\end{itemize}

This kind of local analysis can be applied in both neural and institutional contexts to study whether a given structural configuration is robust under small perturbations of \(E_t\).

\subsection*{Effects of noise and non-smooth landscapes}

In realistic settings:

\begin{itemize}[leftmargin=2em]
    \item the coherence landscape \(C(E, I)\) may be non-smooth or non-convex,
    \item coherence evaluations may be noisy,
    \item and impulsive changes in \(I_t\) can create abrupt landscape deformations.
\end{itemize}

In such cases, $\Delta E$ dynamics may be better modelled as a stochastic process:
\[
    \frac{dE}{dt} = \eta \nabla_E C(E, I) + \xi(t),
\]
with \(\xi(t)\) representing noise. This opens the door to:

\begin{itemize}[leftmargin=2em]
    \item \textbf{Stochastic stability} analysis (e.g.\ using stochastic differential equations and invariant measures),
    \item \textbf{Metastability} and noise-assisted transitions between coherence basins (analogous to simulated annealing),
    \item \textbf{Non-smooth dynamics} where subgradient methods or discrete rule-based updates play the role of \(H_t\).
\end{itemize}

The general UTAM insight remains: whatever the precise form, $\Delta E$ dynamics aim to track coherence gradients (in a broad sense) under noisy, changing conditions, and their qualitative properties (convergence, oscillation, collapse) can be studied with the tools of dynamical systems.

\section{Category-Theoretic UTAM}

To capture the compositional structure of UTAM systems and the distinction between vertical and horizontal layers in an abstract way, it is natural to employ category theory. Here we sketch one possible categorical formulation.

\subsection*{Objects and morphisms}

We begin by considering two intertwined categorical structures:

\begin{itemize}[leftmargin=2em]
    \item a category (or bicategory) encoding \emph{ontological dependence} (vertical UTAM),
    \item and a category encoding \emph{dynamical evolution} (horizontal UTAM).
\end{itemize}

\subsubsection*{Vertical category \(\mathcal{C}_{\mathrm{vert}}\)}

Let \(\mathcal{C}_{\mathrm{vert}}\) be a category whose objects include:

\begin{itemize}[leftmargin=2em]
    \item a distinguished object \(W\) representing will,
    \item objects \(E \in \mathcal{E}(\Will)\) representing structures,
    \item objects \(A \in \mathcal{A}(E)\) representing actions.
\end{itemize}

Morphisms encode ontological dependence:

\begin{itemize}[leftmargin=2em]
    \item \(\phi_E : W \to E\), expressing that \(E\) is a manifestation of \(W\);
    \item \(\psi_{A,E} : E \to A\), expressing that \(A\) is a realization of \(W\) through \(E\).
\end{itemize}

Composition \(W \to E \to A\) corresponds to the UTAM chain \(W \to E \to A\). The essential property is that:

\begin{itemize}[leftmargin=2em]
    \item there are no morphisms \(E \to W\) or \(A \to E\) that invert these arrows in \(\mathcal{C}_{\mathrm{vert}}\);
    \item the vertical structure is thus \emph{asymmetric}, reflecting ontological priority.
\end{itemize}

\subsubsection*{Horizontal category \(\mathcal{C}_{\mathrm{horiz}}\)}

Let \(\mathcal{C}_{\mathrm{horiz}}\) be a category (or more naturally, a category enriched over time-indexed structures) whose objects are time-indexed states of the system:
\[
    X_t := (E_t, I_t, A_t),
\]
and whose morphisms represent dynamical evolution:
\[
    \mu_t : X_t \to X_{t+1},
\]
where \(\mu_t\) encodes:
\[
    A_t = F(E_t, I_t), \quad
    E_{t+1} = G(E_t, A_t, I_t, \Theta), \quad
    I_{t+1} = \text{(environmental update)}.
\]

Composition of morphisms in \(\mathcal{C}_{\mathrm{horiz}}\) corresponds to iterating the dynamics over time.

\subsection*{Vertical–horizontal interaction and compositionality}

To integrate these two layers, we can consider:

\begin{itemize}[leftmargin=2em]
    \item a functor \(U : \mathcal{C}_{\mathrm{horiz}} \to \mathcal{C}_{\mathrm{vert}}\) that ``forgets time'' and projects each state \(X_t\) to its vertical UTAM triple \((W, E_t, A_t)\);
    \item or, dually, a fibred or indexed category structure where, for each object \(W\) in \(\mathcal{C}_{\mathrm{vert}}\), we have a corresponding category of horizontal dynamics over \(\mathcal{E}(W)\).
\end{itemize}

Compositionality then appears in at least two senses:

\begin{enumerate}[leftmargin=2em]
    \item \textbf{System composition.} Two UTAM systems can be combined (e.g.\ an agent and an environment, or two interacting agents) via categorical products, coproducts, or more general monoidal structures. Coherence functionals and $\Delta E$ rules can then be defined on the composite.
    \item \textbf{Layer composition.} Vertical morphisms and horizontal morphisms can be composed in a double-category or 2-category framework, where:
    \begin{itemize}
        \item one dimension represents ontological dependence,
        \item the other represents temporal evolution,
        \item and 2-morphisms encode transformations of dynamics (e.g.\ changing \(F\), \(G\), or \(C\) as part of meta-structural dynamics).
    \end{itemize}
\end{enumerate}

This categorical sketch is intentionally schematic, but it points toward a rigorous formalism in which:

\begin{itemize}[leftmargin=2em]
    \item UTAM systems are objects in a suitable category (or higher category),
    \item their dynamics are morphisms,
    \item and their composition, interconnection, and abstraction can be studied with the standard tools of category theory.
\end{itemize}

In later work, one might develop:

\begin{itemize}[leftmargin=2em]
    \item a precise notion of \emph{UTAM morphism} between systems (structure-preserving maps that respect \(\Will \to E \to A\) and coherence),
    \item monoidal structures capturing parallel and sequential composition,
    \item and adjunctions or limits corresponding to coarse-graining, emergent structures, and effective theories.
\end{itemize}

Such a categorical UTAM would make explicit the intuition that will, structure, and action are not isolated entities but nodes in a network of composable processes, all grounded in the primitive directedness of \(\Will\) and organized through coherence-preserving dynamics.


\chapter{Ontological Status of Will (\texorpdfstring{$W$}{W}} 

UTAM places \(\Will\) at the base of its ontology: an operator of directedness from which structures and actions arise. This raises an inevitable question: \emph{what is the ontological status of \(W\)?} Is it a substantive postulate, a re-labelling of familiar notions (laws, information, powers), or a disguised theological claim? In this chapter we clarify how ONTO$\Sigma$ understands \(W\), and how it relates to truth, laws, and alternative metaphysical pictures.

\section{\texorpdfstring{$W$}{W} as Truthmaker / Ground of Directionality}

\subsection*{Will as ground of directional facts}

In ordinary scientific language, we often speak as if the world simply \emph{behaves} according to certain laws or tendencies: systems minimise free energy, particles follow geodesics, organisms seek homeostasis, agents maximize expected reward. ONTO$\Sigma$ suggests that all of these can be seen as \emph{local expressions of a more primitive fact}: the world exhibits \emph{directedness}. Things do not merely happen; they tend to go somewhere.

We interpret \(\Will\) as the \emph{ground} of directional facts:

\begin{itemize}[leftmargin=2em]
    \item facts of the form ``this system tends to evolve toward $X$'',
    \item ``this configuration is an attractor'',
    \item ``this pattern is stable under perturbation''.
\end{itemize}

In truthmaker terminology, \(\Will\) is not Truth itself but a \emph{truthmaker} (or family of truthmakers) for directional truths. When we say a system \emph{has a tendency}, \emph{seeks a minimum}, or \emph{strives to survive}, we are, in ONTO$\Sigma$ language, describing a particular way in which \(\Will\) has been structured into \(E\) and unfolds as \(A\).

\subsection*{Distinguishing \texorpdfstring{$W$}{W} from Truth}

It might be tempting to identify \(\Will\) with Truth: if all directedness comes from \(\Will\), perhaps truths are just aspects of will. ONTO$\Sigma$ rejects this identification for two reasons:

\begin{enumerate}[leftmargin=2em]
    \item \textbf{Truth is about representation.} Truth, on a broad coherence or correspondence view, concerns the alignment between:
    \begin{itemize}
        \item our \emph{interpretive structures} (cognitive \(E_t^{\mathrm{cog}}\)),
        \item our \emph{statements or models} (symbolic \(A_t^{\mathrm{disc}}\)),
        \item and the actual unfolding of the world.
    \end{itemize}
    \item \textbf{\(\Will\) is about ontological directedness.} \(\Will\) underlies the unfolding itself: it is what makes there \emph{be} a structured evolution for representations to track. It is prior to any particular representation or truth-claim.
\end{enumerate}

Formally, we can say:

\begin{itemize}[leftmargin=2em]
    \item \(\Will\) grounds facts of the form
    \[
        W \to E \to A,
    \]
    i.e.\ that certain structures and actions are possible and have definite directional tendencies.
    \item \emph{Truth} concerns the coherence of \emph{interpretations} with this unfolding:
    \[
        \text{Truth at time } t \quad \approx \quad \mathrm{Coherence}\bigl(
            E_t^{\mathrm{cog}}, A_t^{\mathrm{disc}} \,\mid\, W \to E_t^{\mathrm{world}} \to A_t^{\mathrm{world}}
        \bigr),
    \]
    where \(E_t^{\mathrm{world}}, A_t^{\mathrm{world}}\) denote the actual world-structure and world-actions.
\end{itemize}

Truth, in ONTO$\Sigma$, is thus a \emph{second-order coherence}: a relationship between cognitive structures and the world’s $W \to E \to A$ unfolding. \(\Will\) is the ontological condition that makes there be such an unfolding in the first place.

\section{Alternatives and Objections}

The introduction of \(\Will\) as an ontological primitive invites comparison with other metaphysical options and raises natural objections.

\subsection*{\texorpdfstring{$W$}{W} vs.\ laws of nature}

One alternative is to treat \emph{laws of nature} as the primary explanatory entities:

\begin{itemize}[leftmargin=2em]
    \item In Humean or regularity-based views, laws summarize observed patterns; there is no deeper directedness beyond the fact that certain configurations tend to follow others.
    \item In non-Humean views, laws are primitive modal facts that govern the behaviour of systems.
\end{itemize}

Objection: if we already have laws, do we need \(\Will\)? Perhaps what ONTO$\Sigma$ calls \(\Will\) is just a rebranding of nomic structure.

\subsection*{\texorpdfstring{$W$}{W} vs.\ structural realism}

Structural realism claims that only \emph{relational structure} is ontologically robust: we should take seriously the mathematical structure of theories, not their putative objects. On this view:

\begin{itemize}[leftmargin=2em]
    \item the world is a web of relations and symmetries (a giant \(E\)-like entity),
    \item change and process are captured as transformations of this structure.
\end{itemize}

Objection: if structure is all there is, why postulate a separate \(\Will\)? Are we not simply adding a metaphysical label to \(\mathcal{E}\) and calling it primitive?

\subsection*{Information-only views}

In ``it-from-bit'' or purely informational ontologies, the fundamental furniture of reality consists of:

\begin{itemize}[leftmargin=2em]
    \item bits or qubits,
    \item information flows,
    \item or computational processes.
\end{itemize}

Directedness is then often treated as emergent from information gradients, algorithmic dynamics, or selection processes.

Objection: if \emph{information processing} already explains complex behaviour and adaptation, is a further primitive of \(\Will\) necessary? Could we not simply say that systems follow information-theoretic imperatives (e.g.\ maximize mutual information, minimize free energy) and leave it at that?

\subsection*{Is \texorpdfstring{$W$}{W} necessary, or just a convenient label?}

A more general worry is that \(\Will\) may be \emph{ontologically idle}:

\begin{itemize}[leftmargin=2em]
    \item Everything described by UTAM might be rephrased in terms of laws, structures, and information dynamics.
    \item In that case, \(\Will\) may seem like an unnecessary metaphysical posit, or a poetic shorthand for ``the fact that the world has directional dynamics''.
\end{itemize}

This is compounded by the choice of terminology: ``will'' carries psychological and theological connotations. Skeptics may suspect that ONTO$\Sigma$ is smuggling in a quasi-theological notion under cover of supposedly neutral metaphysics.

\section{Responses and Clarifications}

ONTO$\Sigma$ acknowledges the force of these objections and responds in two steps: by clarifying what \(\Will\) is \emph{not}, and by explaining what work \(\Will\) does that laws, structures, and information do not straightforwardly accomplish.

\subsection*{What \texorpdfstring{$W$}{W} is not}

First, some negative characterizations:

\begin{itemize}[leftmargin=2em]
    \item \textbf{\(\Will\) is not psychological desire.} It is not an inner feeling, intention, or preference of any particular agent. Those are higher-level manifestations of \(\Will\) through cognitive structures.
    \item \textbf{\(\Will\) is not divine will.} ONTO$\Sigma$ does not posit a personal, omniscient subject whose intentions shape reality. \(\Will\) is impersonal: it does not plan, judge, or choose in a moral sense.
    \item \textbf{\(\Will\) is not an extra force or field.} It does not supplement physics with an additional causal term. Rather, it is the \emph{ontological character} of whatever physics describes: that it is inherently directional and structurable.
\end{itemize}

Thus, \(\Will\) is not a theological import or a ghostly addition to science. It is a way of naming the fact that the world can be \emph{organized into coherent, evolving structures} at all.

\subsection*{What \texorpdfstring{$W$}{W} contributes}

Second, what does \(\Will\) add that competing ontologies leave obscure?

\paragraph{Unifying directionality across levels.}

Laws of nature, information dynamics, and structural transformations are typically formulated at specific levels:

\begin{itemize}[leftmargin=2em]
    \item physical (e.g.\ field equations),
    \item biological (e.g.\ fitness landscapes),
    \item cognitive (e.g.\ Bayesian updating),
    \item social (e.g.\ game-theoretic equilibria).
\end{itemize}

Each offers local accounts of why \emph{these} systems behave directionally. ONTO$\Sigma$ claims that there is a \emph{common structural feature} behind them:

\begin{itemize}[leftmargin=2em]
    \item all involve a space of possibilities,
    \item coherence or viability regions in that space,
    \item and dynamics that non-randomly explore and stabilize such regions.
\end{itemize}

\(\Will\) is the name for this cross-level operator of directedness. UTAM’s \(\Will \to \Struct \to \Act\) schema applies to fundamental physics, neural processes, AI systems, and institutions \emph{in the same formal vocabulary}. \(\Will\) thus plays a unifying role: it is the minimal common denominator behind diverse directional phenomena.

\paragraph{Making explicit what laws and information presuppose.}

Both laws-of-nature and information-only views quietly presuppose directedness:

\begin{itemize}[leftmargin=2em]
    \item Laws are formulated as rules about how states evolve: they implicitly assume that there is a \emph{space of states} and a non-random \emph{flow} on it.
    \item Information-theoretic principles (minimizing free energy, compressing data, maximizing mutual information) assume systems that \emph{behave as if} they care about certain quantities.
\end{itemize}

ONTO$\Sigma$ does not deny the utility of these frameworks; it asks: \emph{what makes such formulations appropriate in the first place?} Why does reality lend itself to description in terms of stable laws and information flows that display consistent directional tendencies?

\(\Will\) is proposed as the minimal ontological enrichment that answers this: the world is not just a set of facts; it is a set of facts structured by primitive directedness. Laws and information-theoretic regularities are then local codifications of how \(\Will\) manifests in particular regimes.

\paragraph{Minimal enrichment, not maximal theology.}

Finally, ONTO$\Sigma$ deliberately avoids inflating \(\Will\) into a full-blown metaphysical or theological system:

\begin{itemize}[leftmargin=2em]
    \item It does \emph{not} attribute consciousness, personality, or moral properties to \(\Will\).
    \item It does \emph{not} posit multiple wills, competing deities, or a teleological plan.
    \item It does posit a primitive: that reality is not merely a static configuration but inherently oriented toward structured unfolding.
\end{itemize}

In this sense, \(\Will\) is a \emph{minimal} addition: it adds one primitive---directedness---to the usual roster of objects, properties, relations, and laws. The payoff is that it:

\begin{itemize}[leftmargin=2em]
    \item unifies physical, biological, cognitive, and social directedness under a single ontological operator,
    \item supports the UTAM hierarchy \(W \to E \to A\) across scales,
    \item and provides a grounding for coherence and drift that is not tied to any particular level of description.
\end{itemize}

Whether this enrichment is finally acceptable is an open philosophical question. ONTO$\Sigma$’s claim is modest but precise: if one wants a single, layered framework that:

\begin{itemize}[leftmargin=2em]
    \item captures directedness wherever it appears,
    \item connects it to coherence and drift,
    \item and remains compatible with existing scientific formalisms,
\end{itemize}

then treating \(\Will\) as a primitive operator of directionality is a natural and, arguably, necessary move. The subsequent parts of this work should be read as an extended test of this claim: can \(\Will\), so understood, illuminate rather than obscure the structures and dynamics we actually observe?


\chapter{ONTO$\Sigma$ III Preview: Volume of Will and Consciousness}

ONTO$\Sigma$ IIb has focused on the dynamics of directedness for a given structural space \(\mathcal{E}(\Will)\): how structures \(E_t\) and actions \(A_t\) evolve under impulses \(I_t\), how coherence and drift arise, and how $\Delta E$ processes preserve viability. ONTO$\Sigma$ III will take a further step and ask:

\begin{quote}
How does the \emph{accessible portion} of \(\mathcal{E}(\Will)\) itself change over time, and what role do consciousness, culture, and civilization play in this expansion?
\end{quote}

To gesture toward this next layer, we introduce the notion of a \emph{volume of will} and sketch how consciousness can be seen as a coherence amplifier and a driver of meta-structural growth, with cosmological implications.

\section{Expansion of Accessible \texorpdfstring{$\mathcal{E}(\Will)$}{E(W)}}

\subsection*{Accessible structural region}

Up to now, \(\mathcal{E}(\Will)\) has denoted the \emph{space of all possible structures} that can manifest \(\Will\). In practice, at any given time \(t\), only a subset of this space is \emph{accessible} to a given system or to a given level of organization (e.g.\ a species, a civilization):

\[
    \mathcal{E}_t^{\mathrm{acc}} \subseteq \mathcal{E}(\Will).
\]

This accessible region encodes:

\begin{itemize}[leftmargin=2em]
    \item what forms of organisation can be realized given current physical, biological, technological, and cultural conditions,
    \item what patterns of coherence \(C(E, I)\) can actually be reached via $\Delta E$ dynamics from existing starting points,
    \item and what ranges of actions \(\mathcal{A}(E)\) are practically available.
\end{itemize}

In this sense, even if \(\mathcal{E}(\Will)\) is fixed as an ontological possibility space, its \emph{effective} portion \(\mathcal{E}_t^{\mathrm{acc}}\) can grow, shrink, or reorganize over time.

\subsection*{Volume of will}

To formalize this, we can endow \(\mathcal{E}(\Will)\) with a measure \(\mu\) (or more generally, some size functional) and define the \emph{volume of will} at time \(t\) as

\[
    V_t := \mu\bigl(\mathcal{E}_t^{\mathrm{acc}}\bigr).
\]

Intuitively:

\begin{itemize}[leftmargin=2em]
    \item \(V_t\) measures how much of \(\Will\)'s structural possibility space is currently realized or reachable by the systems under consideration.
    \item In simple settings (e.g.\ an isolated organism), \(V_t\) may be relatively small and slowly varying.
    \item As new forms of organization appear (neuronal assemblies, social institutions, technologies), \(\mathcal{E}_t^{\mathrm{acc}}\) expands and \(V_t\) increases.
\end{itemize}

A central theme of ONTO$\Sigma$ III will be:

\begin{quote}
Consciousness, culture, and civilization are mechanisms by which the accessible region \(\mathcal{E}_t^{\mathrm{acc}}\) expands, increasing the volume \(V_t\) through which \(\Will\) can be coherently expressed.
\end{quote}

\subsection*{Realized vs.\ potential structures}

We can further distinguish:

\begin{itemize}[leftmargin=2em]
    \item \(\mathcal{E}_t^{\mathrm{real}}\): structures actually instantiated at time \(t\),
    \item \(\mathcal{E}_t^{\mathrm{pot}}\): structures that are not instantiated but reachable via available $\Delta E$ and meta-structural dynamics,
\end{itemize}
with
\[
    \mathcal{E}_t^{\mathrm{real}} \subseteq \mathcal{E}_t^{\mathrm{acc}} = 
    \mathcal{E}_t^{\mathrm{real}} \cup \mathcal{E}_t^{\mathrm{pot}}.
\]

The ``volume of will'' can then be decomposed as
\[
    V_t = \mu\bigl(\mathcal{E}_t^{\mathrm{real}}\bigr) 
        + \mu\bigl(\mathcal{E}_t^{\mathrm{pot}}\bigr),
\]
representing both the actually realized structures and the latent structural possibilities opened by existing organization (e.g.\ a culture that makes new technologies thinkable, or a brain that can learn new skills).

ONTO$\Sigma$ III will explore how different forms of consciousness and social organisation change the ratio between realized and potential volume, and how this relates to growth, stagnation, and collapse.

\section{Consciousness as Coherence Amplifier}

Within this expanding space, consciousness is proposed to play a specific functional role: it acts as a \emph{coherence amplifier} and a \emph{meta-structural operator}.

\subsection*{Enlarging the dimensionality and flexibility of \texorpdfstring{$E_t$}{E\_t}}

A non-conscious adaptive system can still implement I$^2$C and $\Delta E$, but its structural changes are constrained by relatively fixed update rules and limited representational capacity. Consciousness appears, in many contexts, to:

\begin{itemize}[leftmargin=2em]
    \item increase the effective dimensionality of \(E_t\): more features, more abstract representations, richer models of self and world;
    \item enable simultaneous consideration of multiple possible actions \(A_t\) and their long-term coherence implications;
    \item support explicit modelling of \emph{its own} $\Delta E$ processes (``how am I changing?''), thus enabling meta-structural updates.
\end{itemize}

Formally, we can think of conscious systems as those for which:

\begin{itemize}[leftmargin=2em]
    \item the structural manifold \(\mathcal{M}\) becomes higher-dimensional,
    \item the coherence functional \(C\) acquires new arguments (e.g.\ anticipated future states, social constraints),
    \item and the update map \(G\) gains dependence on higher-order information (plans, counterfactuals, narratives).
\end{itemize}

This does not make consciousness \emph{mystical}; it positions it as a particular kind of structural complexity that dramatically enlarges \(\mathcal{E}_t^{\mathrm{acc}}\) and reshapes the coherence landscape.

\subsection*{Enhanced \texorpdfstring{$\Delta E$}{ΔE} and meta-structural capabilities}

Conscious agents typically exhibit:

\begin{itemize}[leftmargin=2em]
    \item \textbf{Explicit meta-cognition:} the ability to represent and evaluate their own structures \(E_t\) (beliefs, habits, values) and to initiate deliberate $\Delta E$ operations (learning, unlearning, reframing).
    \item \textbf{Long-range coherence tracking:} the ability to relate current actions \(A_t\) to long-term coherence regimes (life projects, identities, narratives).
    \item \textbf{Cultural coupling:} participation in shared symbolic structures (language, institutions) that make new regions of \(\mathcal{E}(\Will)\) accessible.
\end{itemize}

In terms of UTAM dynamics, consciousness thus:

\begin{itemize}[leftmargin=2em]
    \item enriches \(F\) (I$^2$C) by adding layers of interpretation (metaphor, counterfactual reasoning, moral evaluation),
    \item enriches \(G\) ($\Delta E$) by adding higher-order updates (changing the coherence criteria themselves, not just fitting existing ones),
    \item and participates in meta-structural processes that propagate new coherence regimes across populations.
\end{itemize}

ONTO$\Sigma$ III will explore consciousness not as a mysterious substance but as a phase in the evolution of volitional systems where coherence management becomes reflexive: will manifests as a structure that can intentionally shape its own space of structures.

\section{Cosmological Implications (Sketch)}

Finally, we gesture toward the cosmological horizon of ONTO$\Sigma$: the idea that the universe can be seen as a progressive unfolding of \(\Will\) through increasingly complex structures \(\mathcal{E}(\Will)\).

\subsection*{Universe as volitional unfolding}

At the most speculative level, we can imagine:

\begin{itemize}[leftmargin=2em]
    \item a primordial phase where \(\Will\) is manifested primarily in simple physical structures: fields, particles, early cosmological dynamics;
    \item subsequent phases where \(\mathcal{E}(\Will)\) acquires new layers: chemistry, self-replicating systems, biological evolution;
    \item later phases where nervous systems, brains, and cultures arise, vastly expanding \(\mathcal{E}_t^{\mathrm{acc}}\) and the volume of will.
\end{itemize}

Life, mind, and civilization then appear not as accidents on a passive stage, but as phases in a \emph{volitional expansion}: stages where more and more of \(\mathcal{E}(\Will)\) is explored and stabilized.

\subsection*{Phases of volitional expansion}

Schematically, we might speak of:

\begin{enumerate}[leftmargin=2em]
    \item \textbf{Physical volition:} directedness expressed as regularities in fundamental dynamics and structure formation (stars, galaxies, chemical networks).
    \item \textbf{Biological volition:} directedness expressed as survival, reproduction, and adaptation in organisms implementing local I$^2$C + $\Delta E$.
    \item \textbf{Cognitive volition:} directedness expressed through conscious systems that model themselves and their worlds, amplifying coherence and opening new structural possibilities.
    \item \textbf{Civilizational volition:} directedness expressed through large-scale institutions, technologies, and shared narratives that reshape the very space of accessible structures.
\end{enumerate}

Each phase does not replace the previous one but \emph{builds on it}: higher-level structures remain grounded in lower-level dynamics, even as they carve out new regions of \(\mathcal{E}(\Will)\).

\subsection*{Toward ONTO$\Sigma$ III}

ONTO$\Sigma$ III will develop this sketch in three directions:

\begin{itemize}[leftmargin=2em]
    \item \textbf{Formalization:} defining the volume of will \(V_t\) more precisely and studying how different kinds of systems (biological, cognitive, social) change \(\mathcal{E}_t^{\mathrm{acc}}\).
    \item \textbf{Consciousness:} elaborating the role of consciousness as a coherence amplifier and meta-structural operator, including its vulnerabilities (e.g.\ self-fragmentation, pathologies of reflection).
    \item \textbf{Cosmology:} exploring whether and how existing cosmological, biological, and cultural theories can be reinterpreted as phases in a UTAM-driven expansion of volitional volume.
\end{itemize}

The aim is not to derive a detailed cosmology from first principles, but to show that once \(\Will\) is taken as primitive, there is a natural way to see the history of the universe as a sequence of increasingly rich manifestations of directedness: from the simplest physical laws to the most intricate forms of self-reflective consciousness and collective organization.


% ================================================
% Part V – Implementation, Code, and Experiments
% ================================================

\part{Implementation, Code, and Experiments}

\chapter{Python Prototypes and Simulation Studies}

The preceding parts of this report have developed UTAM at ontological, dynamical, geometric, and applied levels. To bridge these layers with concrete practice, we include a set of Python prototypes and simulation designs that instantiate the I$^2$C and $\Delta E$ schemes in simple environments.

The goal is twofold:

\begin{itemize}[leftmargin=2em]
    \item to demonstrate, in executable form, how coherence \(C_t\), drift \(D_t\), I$^2$C, and $\Delta E$ behave in toy models;
    \item to provide a scaffold for more sophisticated simulations (higher-dimensional structures, complex environments, multi-agent systems) that can be used as experimental testbeds for UTAM.
\end{itemize}

\section{Minimal \texorpdfstring{$\Delta E$}{ΔE} + I$^2$C Prototype (Appendix A)}

\subsection*{Code overview}

Appendix~A contains a minimal Python implementation of a UTAM-style agent interacting with a drifting one-dimensional environment. The core components are:

\begin{itemize}[leftmargin=2em]
    \item \texttt{DriftingEnvironment}: emits impulses \(I_t\) and defines a ``true'' mapping
    \[
        A^{\mathrm{true}}_t = k_t \, I_t,
    \]
    where the parameter \(k_t\) drifts slowly over time.
    \item \texttt{DeltaEAgent}: implements:
    \begin{itemize}
        \item an I$^2$C action-generator:
        \[
            Z_t = \mathrm{Enc}(I_t, E_t), \quad A_t = E_t \, Z_t,
        \]
        with a trivial encoder \(Z_t = I_t\);
        \item a coherence functional \(C_t\) defined as a Gaussian of the squared error between \(A_t\) and \(A_t^{\mathrm{true}}\);
        \item a $\Delta E$ update:
        \[
            E_{t+1} = E_t + \eta \, \frac{\partial C_t}{\partial E_t},
        \]
        representing coherence-gradient ascent.
    \end{itemize}
    \item \texttt{run\_simulation}: runs the coupled dynamics for a given number of steps, logging history.
    \item \texttt{plot\_results} and \texttt{compare\_with\_without\_delta\_e}: produce diagnostic plots and a numerical summary that illustrate the Drift Law and the effect of $\Delta E$.
\end{itemize}

The script is intentionally compact and self-contained. It can be run as a standalone program, producing both textual output and visualisations.

\subsection*{Mapping to UTAM concepts}

The prototype provides a direct, didactic mapping between code variables and UTAM concepts:

\begin{center}
\begin{tabular}{ll}
\toprule
\textbf{UTAM concept} & \textbf{Code element} \\
\midrule
Impulse \(I_t\) & \texttt{I\_t} (from \texttt{env.step()}) \\
Action \(A_t\) & \texttt{A\_t} (from \texttt{agent.act}) \\
True action \(A^{\mathrm{true}}_t\) & \texttt{A\_true\_t} (environment output) \\
Structure \(E_t\) & \texttt{agent.E} (\(E\)-parameter) \\
Environment parameter \(k_t\) & \texttt{env.k\_history} \\
Coherence \(C_t\) & \texttt{C\_t} (from \texttt{coherence(\dots)}) \\
Drift \(D_t\) & \(1 - C_t\) (implicit) \\
I$^2$C map \(F(E_t, I_t)\) & \texttt{interpret} + \texttt{act} methods \\
$\Delta E$ map \(G(E_t, A_t, I_t, \Theta)\) & \texttt{update\_structure} method \\
Adaptation rate \(\eta\) & \texttt{agent.eta} \\
Coherence sensitivity \(\sigma\) & \texttt{agent.sigma} \\
\bottomrule
\end{tabular}
\end{center}

The environment parameter \(k_t\) plays the role of a drifting ``world structure'', while \(E_t\) is the agent's internal structural estimate. $\Delta E$ dynamics track the coherence gradient between these two.

\subsection*{Example runs and plots}

The script includes several visualisations that make the abstract UTAM notions tangible:

\paragraph{Coherence over time.}

The plot of \(C_t\) versus \(t\) shows how coherence behaves as the environment drifts and the agent adapts. Typical patterns:

\begin{itemize}[leftmargin=2em]
    \item An initial transient where \(C_t\) increases as \(E_t\) moves toward the current \(k_t\).
    \item Fluctuations around a relatively high mean \(C_t\) once adaptation stabilizes.
\end{itemize}

This illustrates the balance between drift in \(k_t\) and $\Delta E$-driven adaptation in \(E_t\).

\paragraph{Structure \(E_t\) vs.\ true \(k_t\).}

A second plot overlays:

\begin{itemize}[leftmargin=2em]
    \item the agent's structure history \(\{E_t\}\),
    \item the environment's true parameter trajectory \(\{k_t\}\).
\end{itemize}

Good adaptation manifests as \(E_t \approx k_t\) over time, i.e.\ the agent's internal structure tracks the drifting environment. Misadaptation (e.g.\ low \(\eta\), noisy gradients) shows up as persistent lag or divergence.

\paragraph{Drift Law demonstration: with vs.\ without \texorpdfstring{$\Delta E$}{ΔE}.}

The comparison function runs two simulations:

\begin{itemize}[leftmargin=2em]
    \item \textbf{With $\Delta E$}: \(\eta > 0\), so \(E_t\) adapts continuously.
    \item \textbf{Without $\Delta E$}: \(\eta = 0\), so \(E_t\) is frozen at its initial value.
\end{itemize}

The resulting plots show:

\begin{itemize}[leftmargin=2em]
    \item coherence \(C_t\) remaining relatively high and stable with $\Delta E$,
    \item coherence \(C_t\) decaying under drift when \(E_t\) is frozen, in line with the Drift Law:
    \[
        \frac{d I}{dt} \neq 0, \quad \frac{\partial G}{\partial I} \approx 0 \quad \Rightarrow \quad \frac{d C_t}{dt} \le 0.
    \]
\end{itemize}

A printed summary reports average coherence in both conditions, quantifying the benefit of structural adaptation.

\section{Extended Simulations (Future Work)}

The minimal prototype confirms, in a toy setting, the basic UTAM dynamics: I$^2$C for action generation, $\Delta E$ for structural adaptation, and the Drift Law when $\Delta E$ is disabled. However, genuine scientific leverage will come from richer simulations. Here we outline several natural extensions.

\subsection*{Higher-dimensional structures \texorpdfstring{$E_t$}{E\_t}}

Instead of a single scalar parameter \(E_t\), we can consider:

\begin{itemize}[leftmargin=2em]
    \item vector-valued \(E_t \in \mathbb{R}^n\) representing multi-parameter controllers, neural network weights, or abstract features;
    \item manifold-valued \(E_t \in \mathcal{M}\) representing more complex structural spaces (e.g.\ policies on Lie groups, graph structures).
\end{itemize}

This will allow us to:

\begin{itemize}[leftmargin=2em]
    \item study $\Delta E$ dynamics on non-trivial coherence landscapes,
    \item investigate phenomena such as local optima, metastability, and noise-driven transitions between attractors,
    \item connect more directly with neural and institutional models introduced in earlier parts.
\end{itemize}

\subsection*{More complex environments}

Beyond the drifting linear map, we can embed UTAM agents in:

\begin{itemize}[leftmargin=2em]
    \item \textbf{Gridworlds and discrete MDPs:} where \(I_t\) encodes local observations, \(A_t\) encodes moves, and coherence \(C_t\) tracks both task performance and structural robustness under changing reward/dynamics.
    \item \textbf{Continuous control tasks:} locomotion, manipulation, or navigation environments with time-varying parameters, where \(\Delta E\) updates both model parameters and control policies.
    \item \textbf{Partially observable settings:} where interpretation \(\mathrm{Enc}(I_t, E_t)\) must maintain an internal belief state, making the distinction between I$^2$C (inference + action) and $\Delta E$ (model adaptation) more vivid.
    \item \textbf{Stochastic and adversarial settings:} where drift includes regime shifts, non-stationary opponents, or changing constraints, highlighting the role of coherence monitoring and $\Delta E$ in maintaining viability.
    \end{itemize}

Each of these environments can be used to define UTAM-style benchmarks (as discussed in the AI and robotics chapter) with explicit coherence and drift metrics.

\subsection*{Multiple agents and social UTAM dynamics}

Finally, we can move from single-agent to multi-agent simulations, bringing together the UTAM machinery and the social/institutional interpretations:

\begin{itemize}[leftmargin=2em]
    \item Simulate populations of UTAM agents with individual \(E_t^{(i)}\), interacting through shared environments and emergent institutions.
    \item Define macro-level structures \(E_t^{\mathrm{macro}}\) (e.g.\ norms, markets, governance rules) that constrain and are updated by micro-level actions.
    \item Implement $\Delta E$ at both levels:
    \begin{itemize}
        \item micro-level learning and adaptation,
        \item macro-level institutional updates triggered by coherence signals (e.g.\ aggregate welfare, stability, trust).
    \end{itemize}
\end{itemize}

Such simulations would allow exploration of:

\begin{itemize}[leftmargin=2em]
    \item how individual and collective coherence interact,
    \item how institutional drift and reform arise from local UTAM dynamics,
    \item and how meta-structural interventions (changing update rules, coherence criteria) affect long-term trajectories.
\end{itemize}

Taken together, the minimal prototype and these proposed extensions form the computational backbone of the ONTO$\Sigma$ research programme. They provide concrete platforms where metaphysical claims about \(\Will\), coherence, and drift can be translated into executable models, quantitative predictions, and, ultimately, empirical tests.

\appendix

\chapter{Python $\Delta E$ + I$^2$C Prototype}
\label{appendix:python-prototype}

\section*{Full Code Listing}

\begin{lstlisting}[
  language=Python,
  caption={UTAM DeltaE + I2C Prototype with Visualization},
  label={lst:utam-deltae-i2c},
  breaklines=true,
  breakatwhitespace=true,
  columns=fullflexible
]

"""
UTAM DeltaE + I2C Prototype
---------------------------

This script implements a minimal toy example of:

- I2C (Impulse -> Interpretation -> Coherence) as the action-generator F
- DeltaE (structure update) as the coherence-preserving updater G

Environment:
    - 1D impulses I_t from a drifting linear environment.
    - "True" environment mapping: A_true = k_t * I_t, where k_t drifts slowly over time.

Agent:
    - Internal structure parameter E_t (its estimate of k_t).
    - Action: A_t = E_t * I_t  (simple linear policy).
    - Coherence C_t is high when A_t ~= A_true.
    - DeltaE update: gradient ascent on C_t with respect to E_t.

This is not meant as a realistic model, but as an executable illustration of:
    - Horizontal UTAM dynamics: A_t = F(E_t, I_t), E_{t+1} = G(E_t, A_t, I_t, Theta)
    - Drift: environment changes (k_t drifts) while E_t tries to keep up via DeltaE.

Key demonstration:
    - With eta = 0 (no DeltaE update): coherence decays over time (Drift Law).
    - With eta > 0: agent maintains coherence by adapting structure.
"""

import math
import random
from dataclasses import dataclass, field
from typing import List, Tuple, Optional
import matplotlib.pyplot as plt
import numpy as np


# ----------------------------
# Environment with slow drift
# ----------------------------

@dataclass
class DriftingEnvironment:
    """
    Simple 1D environment:
        A_true_t = k_t * I_t

    where k_t drifts slowly over time.

    Drift rate controls how quickly the environment changes;
    noise_std controls stochastic variation in impulses.
    """
    k0: float = 1.0            # initial "true" slope
    drift_rate: float = 0.001  # magnitude of drift per step
    noise_std: float = 0.5     # noise in impulses
    step_count: int = 0
    k_history: List[float] = field(default_factory=list)

    def step(self) -> Tuple[float, float]:
        """
        Generate one impulse I_t and return (I_t, A_true_t).
        The true mapping slope k_t drifts over time.
        """
        # True slope k_t drifts slowly and smoothly
        self.step_count += 1
        k_t = self.k0 + self.drift_rate * self.step_count
        self.k_history.append(k_t)

        # Impulse I_t drawn from a simple distribution
        I_t = random.gauss(0.0, 1.0)

        # Optionally add some noise to simulate richer conditions
        I_t += random.gauss(0.0, self.noise_std)

        A_true_t = k_t * I_t
        return I_t, A_true_t


# ----------------------------
# Coherence functional
# ----------------------------

def coherence(E_t: float, I_t: float, A_t: float, A_true_t: float,
              sigma: float = 1.0) -> float:
    """
    Coherence C_t in (0, 1], modeled as a Gaussian of the squared error
    between A_t and A_true_t.

        error = A_t - A_true_t
        C_t = exp( - error**2 / (2 * sigma**2) )

    When A_t == A_true_t, coherence = 1.0.
    As error grows, coherence decays towards 0.0.
    """
    error = A_t - A_true_t
    return math.exp(- (error ** 2) / (2.0 * sigma ** 2))


def coherence_gradient_wrt_E(E_t: float, I_t: float, A_t: float, A_true_t: float,
                             sigma: float = 1.0) -> float:
    """
    Compute dC/dE for the coherence function above, assuming:

        A_t = E_t * I_t

    Derivation (informal):
        C = exp(- (A - A_true)^2 / (2*sigma^2))
        dC/dA = C * ( - (A - A_true) / sigma**2 )
        dA/dE = I_t
        => dC/dE = dC/dA * dA/dE
                 = C * ( - (A - A_true) / sigma**2 ) * I_t
    """
    C = coherence(E_t, I_t, A_t, A_true_t, sigma=sigma)
    error = A_t - A_true_t
    dC_dE = C * (-(error) / (sigma ** 2)) * I_t
    return dC_dE


# ----------------------------
# Agent with I2C + DeltaE
# ----------------------------

@dataclass
class DeltaEAgent:
    """
    Minimal DeltaE agent:

    Structure:
        E_t: internal structure parameter (estimate of the true slope k_t)

    I2C (Impulse -> Interpretation -> Coherence):
        - Impulse: I_t
        - Interpretation: here trivial (Z_t = I_t), but could be any encoder.
        - Coherence: select A_t = E_t * I_t (linear policy),
                     implicitly aiming at high C_t.

    DeltaE update:
        E_{t+1} = E_t + eta * dC/dE_t

    where C is the coherence between the agent's action A_t
    and the environment's true action A_true_t.
    """
    E: float = 0.0              # initial structure parameter
    eta: float = 0.1            # learning rate
    sigma: float = 1.0          # coherence sensitivity
    history: List[dict] = field(default_factory=list)
    E_history: List[float] = field(default_factory=list)

    def interpret(self, I_t: float) -> float:
        """
        Interpretation step of I2C.

        Here we keep it trivial: Z_t = I_t. In a richer model,
        this would be an encoding step using E_t.
        """
        Z_t = I_t
        return Z_t

    def act(self, Z_t: float) -> float:
        """
        Coherence step of I2C: produce an action A_t.

        Here we use a simple linear policy:
            A_t = E_t * Z_t
        """
        return self.E * Z_t

    def update_structure(self, I_t: float, A_t: float, A_true_t: float) -> float:
        """
        DeltaE update rule:

            E_{t+1} = E_t + eta * dC/dE_t

        Returns the new E_t.
        """
        grad = coherence_gradient_wrt_E(self.E, I_t, A_t, A_true_t, sigma=self.sigma)
        self.E = self.E + self.eta * grad
        self.E_history.append(self.E)
        return self.E

    def step(self, I_t: float, A_true_t: float, t: int) -> None:
        """
        Perform one full I2C + DeltaE step:

            1. Interpret impulse (I2C: Interpretation)
            2. Act based on interpretation (I2C: Coherence/action)
            3. Compute coherence
            4. Update structure via DeltaE

        Logs the step into `history`.
        """
        # I2C: Interpretation
        Z_t = self.interpret(I_t)

        # I2C: Coherence/action
        A_t = self.act(Z_t)

        # Coherence measure
        C_t = coherence(self.E, I_t, A_t, A_true_t, sigma=self.sigma)

        # DeltaE: structure update
        E_old = self.E
        E_new = self.update_structure(I_t, A_t, A_true_t)

        self.history.append(
            {
                "t": t,
                "I_t": I_t,
                "A_true_t": A_true_t,
                "A_t": A_t,
                "C_t": C_t,
                "E_old": E_old,
                "E_new": E_new,
            }
        )


# ----------------------------
# Simulation
# ----------------------------

def run_simulation(
    steps: int = 1000,
    env: Optional[DriftingEnvironment] = None,
    agent: Optional[DeltaEAgent] = None,
) -> Tuple[DriftingEnvironment, DeltaEAgent]:
    """
    Run a simple simulation of UTAM-style dynamics with drift:

        - Environment drift: k_t slowly changes over time.
        - Agent tries to maintain coherence by updating E_t via DeltaE.

    Returns the final environment and agent objects (with history).
    """
    if env is None:
        env = DriftingEnvironment(k0=1.0, drift_rate=0.002, noise_std=0.3)

    if agent is None:
        agent = DeltaEAgent(E=0.0, eta=0.1, sigma=1.0)

    for t in range(steps):
        I_t, A_true_t = env.step()
        agent.step(I_t, A_true_t, t)

    return env, agent


# ----------------------------
# Visualization
# ----------------------------

def plot_results(env: DriftingEnvironment, agent: DeltaEAgent, 
                 title: str = "UTAM DeltaE Agent: Coherence and Structure Adaptation") -> None:
    """
    Plot coherence over time and compare agent's structure (E_t)
    with true environment parameter (k_t).
    """
    fig, axes = plt.subplots(2, 2, figsize=(12, 8))
    
    # Time steps
    t_steps = list(range(len(agent.history)))
    
    # Plot 1: Coherence over time
    C_vals = [h["C_t"] for h in agent.history]
    axes[0, 0].plot(t_steps, C_vals, 'b-', alpha=0.7, label='Coherence C_t')
    axes[0, 0].axhline(y=np.mean(C_vals), color='r', linestyle='--', 
                       label=f'Mean C = {np.mean(C_vals):.3f}')
    axes[0, 0].set_xlabel('Time step t')
    axes[0, 0].set_ylabel('Coherence C_t')
    axes[0, 0].set_title('Coherence Over Time')
    axes[0, 0].legend()
    axes[0, 0].grid(True, alpha=0.3)
    
    # Plot 2: Structure parameter E_t vs true k_t
    axes[0, 1].plot(t_steps, agent.E_history, 'g-', label='Agent structure E_t')
    axes[0, 1].plot(t_steps, env.k_history[:len(t_steps)], 'r--', 
                    label='True environment k_t')
    axes[0, 1].set_xlabel('Time step t')
    axes[0, 1].set_ylabel('Parameter value')
    axes[0, 1].set_title('Structure Adaptation: E_t vs True k_t')
    axes[0, 1].legend()
    axes[0, 1].grid(True, alpha=0.3)
    
    # Plot 3: Error distribution
    errors = [h["A_t"] - h["A_true_t"] for h in agent.history]
    axes[1, 0].hist(errors, bins=50, density=True, alpha=0.7)
    axes[1, 0].axvline(x=0, color='k', linestyle='--', label='Zero error')
    axes[1, 0].set_xlabel('Action error (A_t - A_true_t)')
    axes[1, 0].set_ylabel('Density')
    axes[1, 0].set_title('Error Distribution')
    axes[1, 0].legend()
    axes[1, 0].grid(True, alpha=0.3)
    
    # Plot 4: Phase space: E_t vs Coherence
    scatter = axes[1, 1].scatter(agent.E_history, C_vals, 
                                 c=t_steps, cmap='viridis', alpha=0.6)
    axes[1, 1].set_xlabel('Structure parameter E_t')
    axes[1, 1].set_ylabel('Coherence C_t')
    axes[1, 1].set_title('Phase Space: E_t vs Coherence (colored by time)')
    plt.colorbar(scatter, ax=axes[1, 1], label='Time step')
    axes[1, 1].grid(True, alpha=0.3)
    
    plt.suptitle(title, fontsize=14, fontweight='bold')
    plt.tight_layout()
    plt.show()


# ----------------------------
# Comparison: With vs Without DeltaE
# ----------------------------

def compare_with_without_delta_e(steps: int = 1000):
    """
    Run two simulations side by side:
    1. Agent with DeltaE update (eta > 0)
    2. Agent without DeltaE update (eta = 0)
    
    Demonstrates the Drift Law: without adaptation, coherence decays.
    """
    # Agent WITH DeltaE
    env1 = DriftingEnvironment(k0=1.0, drift_rate=0.002, noise_std=0.3)
    agent1 = DeltaEAgent(E=0.0, eta=0.1, sigma=1.0)
    env1, agent1 = run_simulation(steps, env1, agent1)
    
    # Agent WITHOUT DeltaE (frozen structure)
    env2 = DriftingEnvironment(k0=1.0, drift_rate=0.002, noise_std=0.3)
    agent2 = DeltaEAgent(E=0.0, eta=0.0, sigma=1.0)  # eta = 0
    env2, agent2 = run_simulation(steps, env2, agent2)
    
    # Plot comparison
    fig, axes = plt.subplots(1, 2, figsize=(12, 4))
    
    # Coherence comparison
    t_steps = list(range(steps))
    axes[0].plot(t_steps, [h["C_t"] for h in agent1.history], 'g-', 
                 label='With DeltaE (eta=0.1)', alpha=0.8)
    axes[0].plot(t_steps, [h["C_t"] for h in agent2.history], 'r-', 
                 label='Without DeltaE (eta=0.0)', alpha=0.8)
    axes[0].set_xlabel('Time step t')
    axes[0].set_ylabel('Coherence C_t')
    axes[0].set_title('Coherence: With vs Without DeltaE Update')
    axes[0].legend()
    axes[0].grid(True, alpha=0.3)
    
    # Structure parameter comparison
    axes[1].plot(t_steps, agent1.E_history, 'g-', label='E_t with DeltaE', alpha=0.8)
    axes[1].plot(t_steps, [0.0] * steps, 'r--', label='E_t without DeltaE (frozen)', alpha=0.8)
    axes[1].plot(t_steps, env1.k_history[:steps], 'k:', label='True k_t', alpha=0.6)
    axes[1].set_xlabel('Time step t')
    axes[1].set_ylabel('Parameter value')
    axes[1].set_title('Structure Parameter Adaptation')
    axes[1].legend()
    axes[1].grid(True, alpha=0.3)
    
    plt.suptitle('UTAM Demonstration: Drift Law vs DeltaE Adaptation', 
                 fontsize=14, fontweight='bold')
    plt.tight_layout()
    plt.show()
    
    # Print summary statistics
    C_with = np.mean([h["C_t"] for h in agent1.history])
    C_without = np.mean([h["C_t"] for h in agent2.history])
    
    print("="*60)
    print("COMPARISON SUMMARY")
    print("="*60)
    print(f"With DeltaE (eta=0.1):    Average coherence = {C_with:.4f}")
    print(f"Without DeltaE (eta=0.0): Average coherence = {C_without:.4f}")
    print(f"DeltaE improvement:       {100*(C_with - C_without)/C_without:.1f}%")
    print("="*60)


# ----------------------------
# Main execution
# ----------------------------

if __name__ == "__main__":
    print("UTAM DeltaE + I2C Prototype")
    print("="*60)
    
    # Run single simulation with DeltaE
    print("\nRunning simulation with DeltaE adaptation...")
    env, agent = run_simulation(steps=1000)
    
    # Simple textual summary
    last = agent.history[-1]
    avg_coherence = sum(h["C_t"] for h in agent.history) / len(agent.history)
    
    print(f"\nFinal structure parameter E_t: {last['E_new']:.4f}")
    print(f"Final true environment k_t: {env.k_history[-1]:.4f}")
    print(f"Average coherence C_t over run: {avg_coherence:.4f}")
    
    # Show a few sample steps
    print("\nSample steps (every 200th step):")
    for h in agent.history[::200]:
        print(
            f"t={h['t']:4d} | I_t={h['I_t']:+.3f} | "
            f"A_true={h['A_true_t']:+.3f} | A={h['A_t']:+.3f} | "
            f"C_t={h['C_t']:.3f} | E={h['E_old']:.3f}->{h['E_new']:.3f}"
        )
    
    # Plot results
    plot_results(env, agent, title="UTAM DeltaE Agent: Coherence and Structure Adaptation")
    
    # Compare with vs without DeltaE
    print("\n" + "="*60)
    print("Demonstrating the Drift Law: With vs Without DeltaE")
    print("="*60)
compare_with_without_delta_e(steps=1000)
\end{lstlisting}


\chapter{Additional Mathematical Details}
\label{appendix:math-details}

This appendix collects several elementary derivations and proof sketches that underlie the examples and claims made in the main text. The emphasis is on simple, illustrative results rather than maximal generality.

\section{Gradient of the Gaussian Coherence Functional}

In the Python prototype (Appendix~\ref{appendix:python-prototype}), we used a scalar coherence functional of the form
\begin{equation}
    C = C(E, I, A, A^{\mathrm{true}}) 
      = \exp\!\left( -\frac{(A - A^{\mathrm{true}})^2}{2\sigma^2} \right),
    \label{eq:gaussian-coherence}
\end{equation}
with
\begin{equation}
    A = E \cdot I,
\end{equation}
where
\begin{itemize}[leftmargin=2em]
    \item \(E \in \mathbb{R}\) is the agent's scalar structure parameter,
    \item \(I \in \mathbb{R}\) is the current impulse from the environment,
    \item \(A^{\mathrm{true}} \in \mathbb{R}\) is the environment's ``true'' action (e.g.\ \(A^{\mathrm{true}} = k I\) for some drifting \(k\)),
    \item \(\sigma > 0\) is a scale parameter controlling sensitivity to error.
\end{itemize}

\subsection*{Derivative with respect to \texorpdfstring{$E$}{E}}

We wish to compute the gradient of \(C\) with respect to \(E\), assuming \(A = E I\). Set
\[
    e := A - A^{\mathrm{true}} = E I - A^{\mathrm{true}}.
\]
Then
\[
    C(E) = \exp\!\left( -\frac{e^2}{2\sigma^2} \right).
\]

By the chain rule,
\[
    \frac{\partial C}{\partial E}
    = \frac{\partial C}{\partial A} \cdot \frac{\partial A}{\partial E}.
\]

First,
\begin{align}
    \frac{\partial C}{\partial A}
    &= \frac{d}{dA} \exp\!\left( -\frac{(A - A^{\mathrm{true}})^2}{2\sigma^2} \right) \\
    &= \exp\!\left( -\frac{(A - A^{\mathrm{true}})^2}{2\sigma^2} \right)
       \cdot \left( -\frac{1}{2\sigma^2} \cdot 2 (A - A^{\mathrm{true}}) \right) \\
    &= C \cdot \left( -\frac{A - A^{\mathrm{true}}}{\sigma^2} \right) \\
    &= C \cdot \left( -\frac{e}{\sigma^2} \right).
\end{align}

Second,
\begin{equation}
    \frac{\partial A}{\partial E} = \frac{\partial (E I)}{\partial E} = I.
\end{equation}

Therefore,
\begin{align}
    \frac{\partial C}{\partial E}
    &= \frac{\partial C}{\partial A} \cdot \frac{\partial A}{\partial E} \\
    &= C \cdot \left( -\frac{e}{\sigma^2} \right) \cdot I \\
    &= C \cdot \left( -\frac{(A - A^{\mathrm{true}})}{\sigma^2} \right) I.
\end{align}

This is precisely the expression implemented in the function \texttt{coherence\_gradient\_wrt\_E} in the code.

\subsection*{Vector generalization}

If we generalize to a vector-valued structure \(E \in \mathbb{R}^n\) and a linear policy of the form
\[
    A = \langle E, \phi(I) \rangle = E^\top \phi(I),
\]
with a feature map \(\phi : \mathbb{R}^d \to \mathbb{R}^n\), the same derivation yields
\begin{equation}
    \nabla_E C
    = C \cdot \left( -\frac{(A - A^{\mathrm{true}})}{\sigma^2} \right) \phi(I).
\end{equation}
This is the natural multi-dimensional analogue of the scalar gradient above and can be used in higher-dimensional $\Delta E$ updates.

\section{Monotone Ascent of Coherence under Exact Gradient Updates}

We now consider a more abstract setting where coherence is a smooth function of structure:
\[
    C : \mathcal{M} \to \mathbb{R}, \quad E \mapsto C(E),
\]
with \(0 \le C(E) \le 1\) for all \(E \in \mathcal{M}\), and we perform gradient ascent on \(C\) with respect to \(E\).

For simplicity, let \(\mathcal{M} = \mathbb{R}^n\) and fix impulses \(I\) and any dependence on \(A\): write simply \(C(E) := C(E, I)\). Consider the discrete-time update:
\begin{equation}
    E_{t+1} = E_t + \eta \nabla C(E_t),
    \label{eq:gradient-ascent}
\end{equation}
with learning rate \(\eta > 0\).

\subsection*{L-smooth coherence functions}

Assume that \(C\) is \emph{L-smooth}, i.e.\ its gradient is Lipschitz continuous with constant \(L > 0\):
\begin{equation}
    \|\nabla C(E) - \nabla C(F)\| \le L \|E - F\| \quad \forall E, F \in \mathbb{R}^n.
\end{equation}

Standard results in optimization theory give the following descent-type inequality for smooth functions, which we apply in an ascent context.

\subsection*{One-step improvement inequality}

For any \(E, Y \in \mathbb{R}^n\),
\begin{equation}
    C(Y) \ge C(E) + \nabla C(E)^\top (Y - E) - \frac{L}{2} \|Y - E\|^2.
    \label{eq:l-smooth-inequality}
\end{equation}

Now let \(Y = E + \eta \nabla C(E)\), corresponding to a single gradient ascent step. Then
\[
    Y - E = \eta \nabla C(E),
\]
and substituting into \eqref{eq:l-smooth-inequality} gives
\begin{align}
    C(E + \eta \nabla C(E))
    &\ge C(E) + \nabla C(E)^\top (\eta \nabla C(E))
       - \frac{L}{2} \|\eta \nabla C(E)\|^2 \\
    &= C(E) + \eta \|\nabla C(E)\|^2 - \frac{L\eta^2}{2} \|\nabla C(E)\|^2 \\
    &= C(E) + \left(\eta - \frac{L\eta^2}{2}\right) \|\nabla C(E)\|^2.
\end{align}

\subsection*{Theorem: monotonic ascent for small step size}

\textbf{Theorem B.1 (Monotone ascent in L-smooth case).}
Suppose \(C : \mathbb{R}^n \to \mathbb{R}\) is L-smooth and we update \(E_t\) according to \eqref{eq:gradient-ascent}. If the step size satisfies
\[
    0 < \eta < \frac{2}{L},
\]
then the sequence \(\{C(E_t)\}_{t \ge 0}\) is non-decreasing:
\[
    C(E_{t+1}) \ge C(E_t) \quad \forall t.
\]

\emph{Proof sketch.}
Apply the one-step inequality with \(E = E_t\). For each \(t\),
\[
    C(E_{t+1})
    = C(E_t + \eta \nabla C(E_t))
    \ge C(E_t) + \left(\eta - \frac{L\eta^2}{2}\right) \|\nabla C(E_t)\|^2.
\]
If \(0 < \eta < 2/L\), then
\[
    \eta - \frac{L\eta^2}{2} > 0,
\]
and hence
\[
    C(E_{t+1}) - C(E_t) \ge 0.
\]
\hfill\(\square\)

\subsection*{Boundedness and convergence of coherence values}

Since by construction \(0 \le C(E) \le 1\) for all \(E\), the non-decreasing sequence \(C(E_t)\) is bounded above and therefore converges:
\[
    C(E_t) \to C_\infty \in [0,1] \quad \text{as } t \to \infty.
\]

One can further show (under standard assumptions) that the gradient norms tend to zero along subsequences: any limit point \(E^\ast\) of \(\{E_t\}\) satisfies \(\nabla C(E^\ast) = 0\), i.e.\ is a critical point of \(C\). This is the content of classical convergence results for gradient methods on smooth functions.

In UTAM terms, these facts justify the intuition that, in static environments with well-behaved coherence landscapes and sufficiently small step sizes, $\Delta E$ dynamics will monotonically increase coherence and approach a locally coherent regime.

\section{Boundedness of Trajectories under Simple Assumptions}

We briefly sketch a simple boundedness result for $\Delta E$ dynamics that is often implicitly assumed in toy models.

\subsection*{Setup}

Assume:
\begin{itemize}[leftmargin=2em]
    \item The structural manifold is \(\mathcal{M} = \mathbb{R}^n\).
    \item The coherence function \(C(E)\) is continuous and satisfies \(0 \le C(E) \le 1\).
    \item The gradient is bounded:
    \[
        \|\nabla C(E)\| \le G_{\max} \quad \forall E.
    \]
\end{itemize}

Consider updates \(E_{t+1} = E_t + \eta \nabla C(E_t)\) with constant \(\eta > 0\).

\subsection*{Simple bound}

We can bound the distance from the initial structure \(E_0\) after \(T\) steps:
\begin{align}
    \|E_T - E_0\|
    &= \left\| \sum_{t=0}^{T-1} (E_{t+1} - E_t) \right\| \\
    &= \left\| \sum_{t=0}^{T-1} \eta \nabla C(E_t) \right\| \\
    &\le \sum_{t=0}^{T-1} \eta \|\nabla C(E_t)\| \\
    &\le \sum_{t=0}^{T-1} \eta G_{\max} \\
    &= T \eta G_{\max}.
\end{align}

Thus, for any finite time horizon \(T\),
\begin{equation}
    \|E_T - E_0\| \le T \eta G_{\max}.
\end{equation}
If in addition we assume the system runs under a finite time budget or that \(\eta\) decays over time (e.g.\ \(\eta_t \to 0\)), one can obtain stronger bounds showing that \(E_t\) remains within a bounded region around the initial state.

In more realistic models, one often restricts attention to a compact feasible set \(\mathcal{K} \subset \mathbb{R}^n\) and projects \(E_{t+1}\) back onto \(\mathcal{K}\) if necessary. This yields standard projected-gradient dynamics, whose boundedness and convergence properties are well-studied.

\section{Coupled Dynamics with Drift (Sketch)}

Finally, we briefly comment on the coupled system
\[
    \frac{dE}{dt} = \eta \nabla_E C(E, I), \qquad
    \frac{dI}{dt} = f(I, \text{environment}),
\]
introduced in the main text as a stylised model of adaptation under drift.

A full analysis requires specifying \(f\) and \(C\); here we simply note two generic phenomena:

\begin{itemize}[leftmargin=2em]
    \item If \(I(t)\) varies slowly relative to the timescale of $\Delta E$ adaptation (small \(\|\dot{I}(t)\|\)), then under suitable regularity conditions one can approximate \(E(t)\) as tracking a moving local maximizer of the ``instantaneous'' coherence function \(C(\cdot, I(t))\).
    \item If \(I(t)\) changes too quickly, or if the coherence landscape \(C(\cdot, I)\) undergoes bifurcations (e.g.\ local maxima appearing/disappearing), then tracking can fail, leading to episodes of sharp coherence loss---a continuous analogue of the Drift Law.
\end{itemize}

In the simplest quasi-static regime, one can use tools from singular perturbation theory and adiabatic tracking to formalize these intuitions; in more complex regimes, one must rely on numerical simulation and qualitative dynamical-systems analysis. Developing such results more fully is a natural direction for future technical work within the ONTO$\Sigma$ programme.


\chapter{Glossary}
\label{appendix:glossary}

This glossary provides quick-reference definitions of the core concepts and symbols used throughout the ONTO$\Sigma$ framework.

\begin{description}[leftmargin=2.5cm,style=nextline]

    \item[\textbf{Will} (\(W\))] 
    The primitive ontological operator of directedness. Not psychological desire or divine intention, but the fundamental ``vector of being'' that makes structured, goal-like unfolding possible at all. All structures and actions are manifestations of \(\Will\).

    \item[\textbf{Structure} (\(E\))] 
    A concrete organization or pattern through which \(\Will\) is embodied. Includes physical configurations, generative models, controllers, organisms, cognitive schemas, institutions, etc. In UTAM, \(E\) is always a \emph{structuring of \(\Will\)}, not an independent substance.

    \item[\(\mathcal{E}(\Will)\) (Structure space)] 
    The space of all structures that can manifest \(\Will\). Each \(E \in \mathcal{E}(\Will)\) is a possible way in which directedness can be organized and stabilized.

    \item[\textbf{Action} (\(A\))] 
    The unfolding of a structure \(E\) over time. A concrete, observable realization of \(\Will\) through a particular structure. Includes physical motions, control outputs, behaviours, institutional decisions, etc.

    \item[\(\mathcal{A}(E)\) (Action space)] 
    The set of possible actions realizable through a given structure \(E\). For different \(E\), the corresponding \(\mathcal{A}(E)\) may differ radically.

    \item[\textbf{UTAM axiom} (\(W \to E \to A\))] 
    The core ontological schema: \emph{Primacy of Will}, \emph{Mediation by Structure}, \emph{Actualization as Action}. No \(E\) without \(W\); no \(A\) without \(E\). The arrow \(\to\) here expresses ontological \emph{priority}, not temporal succession.

    \item[\textbf{Vertical UTAM}] 
    The atemporal, ontological reading of \(W \to E \to A\). Encodes dependence relations (will $\Rightarrow$ structure $\Rightarrow$ action) without reference to time or feedback.

    \item[\textbf{Horizontal UTAM}] 
    The temporal, dynamical layer. Describes how structures and actions evolve in time under impulses and adaptation:
    \[
        A_t = F(E_t, I_t), \qquad
        E_{t+1} = G(E_t, A_t, I_t, \Theta).
    \]
    Feedback operates entirely within \(\mathcal{E}(\Will)\), without creating or destroying \(\Will\).

    \item[\textbf{Impulse} (\(I_t\))] 
    An input or disturbance from the environment at time \(t\). Includes sensory data, tasks, shocks, social signals, physiological changes, etc. The ``raw'' material interpreted by the system.

    \item[\(\mathbf{F}\) (Action-generator)] 
    The horizontal map \(F(E_t, I_t)\) that produces an action \(A_t\) from the current structure and impulse. Operationally implemented as the I$^2$C cycle (Impulse $\to$ Interpretation $\to$ Coherence/action).

    \item[\(\mathbf{G}\) (Structure-updater)] 
    The horizontal map \(G(E_t, A_t, I_t, \Theta)\) that updates structure in light of actions and impulses. Operationally implemented as $\Delta E$ (coherence-preserving structural adaptation).

    \item[\(\Theta\) (Adaptation parameters)] 
    Parameters controlling the behaviour of \(G\): learning rates, plasticity coefficients, regularisation strengths, memory time-scales, etc.

    \item[\textbf{Coherence} (\(C_t\))] 
    A scalar (or vector) measure in \([0,1]\) of how well structure \(E_t\) is aligned with impulses \(I_t\) and actions \(A_t\) at time \(t\). Can be instantiated via predictive accuracy, stability margins, viability, narrative integrity, phenomenological ``rightness'', etc.

    \item[\textbf{Drift} (\(D_t\))] 
    The complement of coherence at time \(t\):
    \[
        D_t = 1 - C_t.
    \]
    Measures misalignment between structure and environment, or between structure and its own actions. Drift increases when structures fail to adapt to changing impulses.

    \item[\textbf{Drift Law}] 
    A qualitative law of horizontal UTAM: in a non-stationary environment (\(\dot{I} \neq 0\)), if the structure-updater \(G\) is effectively insensitive to impulses (\(\partial G / \partial I \approx 0\), or frozen \(\Theta\)), then coherence tends to decrease and drift tends to increase:
    \[
        \frac{d C_t}{dt} \le 0, \qquad \frac{d D_t}{dt} \ge 0.
    \]

    \item[\textbf{I$^2$C cycle}] 
    The Impulse $\to$ Interpretation $\to$ Coherence pattern implementing \(F\):
    \[
        I_t \;\to\; Z_t = \mathrm{Enc}(I_t, E_t)
        \;\to\; A_t = \arg\max_{A} C(Z_t, A, E_t).
    \]
    A general schema for how an adaptive system turns raw impulses into coherent actions via internal representation.

    \item[\(\Delta E\) (\textbf{Delta-E} update)] 
    The principle of structure adaptation aimed at preserving or increasing coherence under drift. Canonical gradient form:
    \[
        E_{t+1} = E_t + \eta \nabla_E C(E_t, I_t, A_t),
    \]
    but conceptually any update rule that tends to increase coherence (in expectation) counts as $\Delta E$ (gradient-based, Bayesian, Hebbian, institutional reform, etc.).

    \item[\(\eta\) (Learning rate / adaptation speed)] 
    Step-size parameter in $\Delta E$ updates. Controls how quickly structure moves along coherence gradients. Too small $\Rightarrow$ slow adaptation; too large $\Rightarrow$ instability or overshoot.

    \item[\textbf{Volitional phase space}] 
    The geometric space in which UTAM dynamics are visualized:
    \[
        \mathcal{V} \approx \mathcal{E}(\Will) \times \mathcal{I},
    \]
    where \(\mathcal{E}(\Will)\) is the structural manifold and \(\mathcal{I}\) the space of impulses. Points represent \((E, I)\) pairs; trajectories represent adaptive evolution under I$^2$C + $\Delta E$.

    \item[\textbf{Structural manifold} (\(\mathcal{M}\))] 
    A geometric model of the structure space \(\mathcal{E}(\Will)\). Often taken as a differentiable manifold where gradient-based $\Delta E$ dynamics can be defined.

    \item[\textbf{Coherence landscape}] 
    The scalar field
    \[
        C : \mathcal{E}(\Will) \times \mathcal{I} \to [0,1],
    \]
    viewed as a landscape of peaks (high coherence) and valleys (low coherence) over the volitional phase space. $\Delta E$ dynamics correspond to moving uphill on this landscape.

    \item[\textbf{Attractor basin}] 
    A region of the structural manifold where $\Delta E$ dynamics tend to converge to a stable coherence regime (local maximum of \(C\)). Psychologically associated with ``flow states'', habits, or stable identities; socially with institutional equilibria.

    \item[\textbf{Meta-structural dynamics}] 
    Processes that change not only the current structure \(E_t\) but the \emph{accessible region} and update rules themselves. Includes conscious reflection, cultural innovation, technological change, and institutional redesign that expand or reshape \(\mathcal{E}_t^{\mathrm{acc}}\).

    \item[\(\mathcal{E}_t^{\mathrm{acc}}\) (Accessible structure region)] 
    The subset of \(\mathcal{E}(\Will)\) that is reachable or realizable at time \(t\) given current physical, biological, cognitive, and cultural constraints.

    \item[\(\mathcal{E}_t^{\mathrm{real}}\) (Realized structures)] 
    The subset of \(\mathcal{E}_t^{\mathrm{acc}}\) actually instantiated at time \(t\) in concrete systems (organisms, agents, institutions, technologies).

    \item[\textbf{Volume of will} (\(V_t\))] 
    A measure of the size of the accessible region:
    \[
        V_t := \mu\bigl(\mathcal{E}_t^{\mathrm{acc}}\bigr),
    \]
    for some measure \(\mu\) on \(\mathcal{E}(\Will)\). Intuitively, how much of \(\Will\)'s structural possibility space is currently open and explorable.

    \item[\textbf{Self (Self\(_t\))}] 
    In ONTO$\Sigma$, the self at time \(t\) is understood as a coherence structure tying together \(\Will\), attention, interpretation, memory, and action. A relatively stable attractor in volitional phase space for a given organism or person.

    \item[\textbf{Meaning (long-term coherence)}] 
    A long-term coherence regime: a trajectory in volitional phase space where \(C_t\) remains high across a wide range of impulses, often associated with stable identities, life projects, and narrative integrity.

    \item[\textbf{Institutional structure} (\(E_t^{\mathrm{macro}}\))] 
    Collective-level structures (laws, norms, infrastructures, organizations) that mediate \(\Will\) at social and civilizational scales. They shape and constrain the actions of many individual agents.

    \item[\textbf{Civilizational drift}] 
    Long-term increase in misalignment between large-scale structures \(E_t^{\mathrm{civ}}\) (institutions, infrastructures, dominant narratives) and underlying conditions (ecological, technological, demographic), often due to failed or captured $\Delta E$ at macro-scale.

\end{description}


\chapter{ONTO$\Sigma$ References}
\label{appendix:ontos-references}

This appendix summarizes the core ONTO$\Sigma$ texts that form the conceptual and technical background for the present report. Each item includes a brief overview and a placeholder link for citation and distribution.\footnote{Replace the placeholder URLs with actual links or DOIs when available.}

\section*{ONTO$\Sigma$ I: ``Will as an Ontological Operator''}

\subsection*{Summary}

ONTO$\Sigma$ I introduces the central thesis that \emph{will} (\(\Will\)) is an \emph{ontological operator of directedness}---a primitive, impersonal ``vector of being'' that grounds the fact that processes in the world are not merely causal but \emph{directed}. The text:

\begin{itemize}[leftmargin=2em]
    \item motivates the move from purely mechanistic or law-based ontologies to a volitional ontology where directedness is fundamental rather than derivative;
    \item carefully distinguishes \(\Will\) from:
    \begin{itemize}
        \item subjective psychological desire or choice,
        \item divine or theological will,
        \item and ordinary physical forces or fields;
    \end{itemize}
    \item formulates the UTAM axiom in its purely vertical form:
    \[
        \Will \to \Struct \to \Act,
    \]
    where \(\Struct\) denotes structure and \(\Act\) denotes action;
    \item discusses historical precedents (Schopenhauer, Bergson, Heidegger, process metaphysics) and clarifies how ONTO$\Sigma$ agrees with and diverges from them, especially in treating will as \emph{structured and purposive} rather than blind striving;
    \item argues that many recurrent technical patterns observed in the \emph{Synthetic Conscience} sequence (coherence measures, thermostat-like dynamics, $\Delta E$ patterns) are local, technical shadows of a deeper ontological layer—\(\Will\) as the primordial operator of directedness.
\end{itemize}

\noindent\textbf{Citation (working):}
\begin{quote}
M.\ Barzenkov, \emph{ONTO$\Sigma$ I: ``Will'' as an Ontological Operator}, preprint, 2025.
\end{quote}

\noindent\textbf{Link (to be updated):}
\begin{quote}
\url{https://example.org/ontos/ontos-I-will-as-ontological-operator.pdf}
\end{quote}

\vspace{1em}

\section*{ONTO$\Sigma$ II: ``Unified Theory of Adaptive Meaning (UTAM)''}

\subsection*{Summary}

ONTO$\Sigma$ II develops the \emph{Unified Theory of Adaptive Meaning} (UTAM), extending the vertical axiom \(\Will \to \Struct \to \Act\) into a framework capable of describing adaptive systems across physical, biological, cognitive, and social domains. The text:

\begin{itemize}[leftmargin=2em]
    \item refines the UTAM axiom into three tightly coupled postulates:
    \begin{enumerate}
        \item \emph{Primacy of Will (\(\Will\))}: there is a fundamental operator of directedness;
        \item \emph{Mediation by Structure (\(\Struct\))}: will never appears ``naked'' but only through organized structures (physical, informational, social);
        \item \emph{Actualization as Action (\(\Act\))}: actions are the temporal unfolding of structures, and thus indirect manifestations of \(\Will\);
    \end{enumerate}
    \item introduces the notions of \emph{coherence} and \emph{drift} as measures of alignment and misalignment between \(\Will\), structure, and action under changing conditions;
    \item distinguishes \emph{vertical UTAM} (ontological priority) from \emph{horizontal UTAM} (temporal dynamics and feedback);
    \item sketches the Drift Law: in non-stationary environments, structures that fail to update in response to new impulses inevitably fall out of coherence;
    \item motivates the need for explicit dynamical schemes (later elaborated as I$^2$C and $\Delta E$ in ONTO$\Sigma$ IIb) to connect metaphysical insight with operational models.
\end{itemize}

ONTO$\Sigma$ II thus serves as the bridge between the purely ontological claims of ONTO$\Sigma$ I and the fully dynamical, algorithmic, and geometric elaborations developed in ONTO$\Sigma$ IIb (this report).

\noindent\textbf{Citation (working):}
\begin{quote}
M.\ Barzenkov, \emph{ONTO$\Sigma$ II: Unified Theory of Adaptive Meaning (UTAM)}, preprint, 2025.
\end{quote}

\noindent\textbf{Link (to be updated):}
\begin{quote}
\url{https://example.org/ontos/ontos-II-unified-theory-of-adaptive-meaning.pdf}
\end{quote}

\vspace{1em}

\section*{ONTO$\Sigma$ III: ``The Volume of Will and the Role of Consciousness'' (in progress)}

\subsection*{Summary}

ONTO$\Sigma$ III, currently in draft form, extends UTAM in a cosmological and phenomenological direction. It introduces the notion of the \emph{volume of will} and investigates the role of consciousness and culture in expanding the accessible region of the structural space \(\mathcal{E}(\Will)\). The main themes include:

\begin{itemize}[leftmargin=2em]
    \item distinguishing the full possibility space \(\mathcal{E}(\Will)\) from the \emph{accessible} subset \(\mathcal{E}_t^{\mathrm{acc}}\) at a given time and scale;
    \item defining the \emph{volume of will} \(V_t = \mu(\mathcal{E}_t^{\mathrm{acc}})\) as a measure of how much of \(\Will\)'s structural potential is currently realizable or reachable;
    \item interpreting consciousness as a \emph{coherence amplifier} and a \emph{meta-structural operator} that:
    \begin{itemize}
        \item increases the effective dimensionality and flexibility of structures \(E_t\),
        \item enables explicit meta-cognition and deliberate $\Delta E$,
        \item couples individual adaptation to cultural and institutional dynamics;
    \end{itemize}
    \item sketching a cosmological narrative in which the universe is seen as a progressive unfolding of \(\Will\) through increasingly complex structures:
    \begin{enumerate}
        \item physical volition (basic laws and structure formation),
        \item biological volition (life, adaptation, survival),
        \item cognitive volition (mind, modelling, self-awareness),
        \item civilizational volition (institutions, technology, collective reflection).
    \end{enumerate}
\end{itemize}

Within this arc, ONTO$\Sigma$ III asks how consciousness and civilization participate in expanding \(V_t\), and what ethical and practical constraints might apply to large-scale meta-structural interventions.

\noindent\textbf{Citation (working):}
\begin{quote}
M.\ Barzenkov, \emph{ONTO$\Sigma$ III: The Volume of Will and the Role of Consciousness}, draft manuscript, 2025.
\end{quote}

\noindent\textbf{Link (to be updated):}
\begin{quote}
\url{https://example.org/ontos/ontos-III-volume-of-will-and-consciousness.pdf}
\end{quote}


\nocite{*}
\bibliographystyle{unsrt}
\bibliography{references}
\clearpage

\end{document}
